\begin{chapterbox}
    \chapter{Der größte Seeräuber der Nordmeere (WIP)}
    \label{wip-der-groesste-seeraeuber-der-nordmeere}
    \az{Jahr 66}
    
    Eine Freibeuterin, ein Meeresgeist, ein Schiffsbauer und ein Meisterhandwerker vereinen ihre Fähigkeiten. Nur gemeinsam können sie den Zwangswerften der \gend{Seekrieger*innen}{Seekrieger} entgehen und vielleicht die Karte zum größten Schatz der Nordmeere ergattern.
\end{chapterbox}



Diese Fan-Geschichte ist noch in Arbeit. Ich teile sie bereits hier, damit sie für mich nicht verloren geht, falls mein Laptop spontan verschwinden sollte, und damit sie für euch nicht völlig verloren geht, falls ich spontan verschwinden sollte. Fühlt euch in letzterem Fall vollkommen frei, hierrin zu schmökern und nach Lust und Laune weiter daran zu werkeln. Aber solange ich ein aktiver Andori bleibe, lest bitte nicht zu genau nach, was ich geplant haben könnte. LG BBB :D



\newpage
\section{Cybelle, die Freibeuterin}
\az{Jahr 66}

Das dunkle Schiffssegel bauschte sich auf im salzigen Meereswind. Darauf abgebildet war eine schaurige schwarze Seeschlange. Das Zeichen der Schwarzen Kogge, des gefürchteten Piratenschiffs von Kapitän Callem und dessen legendärer Besatzung.

Dieses Segel hier gehörte natürlich nicht der Schwarzen Kogge selbst – jenes Schiff war seit vielen Jahrdutzenden verflucht und verschollen – sondern zur \textsc{Perlenechse}, dem kleinen Schiff der Freibeuterin Cybelle, welche soeben sehr selbstgerecht hinter dem Steuer stand und Befehle gegen den Orkan brüllte: "Praeda, bereite die Enterhaken vor! Minerma, halte deine Waffen bereit! Skyld, bemanne die Bugballista! Corsaria Algenbart, spähe nach Kreaturen aus Arkterons Schlund!"

"Kapitänin, ich erspähe vor allem, dass der Abstand zur \textsc{Capella} immer größer wird!", brüllte Corsaria zurück, den Blick nie von ihrem Fernrohr lassend.

Da hatte sie recht, musste Cybelle zähneknirschend zugeben. Die \textsc{Capella}, das erheblich größere Schiff der Seehändlerin Mondrianne, legte weitere Knoten zu. Bald würde es nicht mehr zu erkennen sein durch Gischt und Sturmwolken hindurch. Wieder einmal würde die Händlerin es unterlassen, den angemessenen Tribut für das Durchqueren der Region um die Echsenfinger-Klippen zu berappen. Doch so rasch würde Cybelle diesmal nicht aufgeben. 

"Praeda, zum Steuer!", befahl sie. Der kleine Praeda bellte bejahend und schlitterte über das glitschige Deck zum Steuer, damit Cybelle es verlassen konnte. Cybelle nahm ihm im vorbei Stolpern einen Enterhaken ab, klammerte sich mit ihrer einen Hand an die Reling und wirbelte mit der anderen den Enterhaken am langen dicken Seil um sich herum. Einmal losgelassen, vollführte das Metallteil einen eleganten Bogen durch den dichten Nebel hindurch und verhakte sich mit dem Heck der \textsc{Capella}. Geschickt knotete Cybelle das andere Ende an eigens dafür gedachte Holzpfeiler auf ihrer Seite. Es knarzte gefährlich, als das Seil sich mit einem Ruck spannte. Die \textsc{Perlenechse} schnellte einmal, zweimal auf die \textsc{Capella} zu. 

An Bord der \textsc{Capella} wagte ein junger Matrose, sich über die Reling zu beugen und mit einem gekrümmten Messer am Seil zu schneiden. Beim Anblick seines erfolglosen Vorhabens konnte Cybelle ein Glucksen kaum unterdrücken. Ihre teuren Seile waren gedreht aus hochstabilen tulgorischen Pflanzenfasern, aus denen Tulgori feuer- und reißfeste Kleidungsstücke und sogar Fischernetze herstellten, die den größten Säbelfischen widerstanden. Ein mickriges Messer würde sie nicht zerteilen können. Ganz zu schweigen davon, dass inzwischen bereits ihr zweiter Enterhaken sein Ziel gefunden hatte. 

Cybelle stieß einen heroischen Schlachtruf aus, klemmte ihren Säbel zwischen ihre Zähne und hangelte sich geübt am glitschigen Seil entlang auf das Händlerschiff zu. Hinter ihr tat Minerma es ihr gleich. Der junge Matrose am gegnerischen Bord wich ängstlich zurück.

Unter ihnen blubberte es aufgeregt im stürmischen Ozean. Golden glitzerndes Haar brach durchs Wasser. Unzweifelhaft waren dies \gend{Nix*en}{Nixen}, welche sich für das Geschehen an der Wasseroberfläche interessierten. Doch hielten die \gend{Nix*en}{Nixen} gehörig Abstand zur \textsc{Capella} und der kleinen Nussschale, die sich daran gehängt hatte.

Keuchend schwang Cybelle sich über die Reling, fuchtelte mit ihrem Säbel herum und sondierte die Lage an Bord des Handelsschiffs. Hier nahm der junge Matrose seine Beine in die Hand, dort kämpfte eine Matrosin mit einem im Sturmwind flatternden Laken, am Steuer knurrte die altbekannte Kapitänin Mondrianne einige kreative Flüche – allem Anschein nach beschützten keine \gend{Seekrieger*innen}{Seekrieger} die Ladung, das Schiff verfügte über keine Ballisten und niemand war bereit, ernsthaft Widerstand zu leisten. Perfekt.

Cybelle umrundete den kleinen Honignussbaum, der unter dem Hauptmast in einer riesigen Kiste Erde stand. Komplett sinnbefreit, einen Baum auf einem Schiff wachsen zu lassen versuchen, aber was man nicht alles für seine Religion tat.

"Nur keine Angst, wir sind nicht hier, um Eure Honignüsse zu stehlen oder Euch um all Euer teuer erfeilschtes Hab und Gut zu bringen", rief Cybelle, Selbstsicherheit projizierend. Mann, in solchen Momenten fühlte sie sich erst richtig lebendig! "Wir sind nur hier, um angemessenen Zoll für das Durchqueren der Echsenfinger-Klippen einzutreiben. Plus eine Entschädigung für unseren Mehraufwand entstanden durch Euren sinnlosen Fluchtversuch. Und tut nicht so, als würden die reichen \gend{Schnösel*innen}{Schnösel} von Werftheim nicht ebenfalls Tribut für den Warentransport durch von ihnen patrouillierte Gewässer verlangen."

Während sie sprach, näherte sie sich langsam der weiterhin mit dem Steuer ringenden Kapitänin Mondrianne. Diese drehte ihr tätowiertes Gesicht zu Cybelle und spuckte aus.

Ungerührt fuhr Cybelle fort: "Eure Spende hier wird sogar dem gemeinen Volk der Nebelinseln zugutekommen, statt die prallen Taschen der reichsten \gend{Schiffsbaumeister*innen}{Schiffsbaumeister} weiter zu füllen. Das verspreche ich euch in meinem Namen. Und der hat Gewicht, schließlich bin ich die größte Freibeuterin der Nordmeere, Cy..."

"Seeräuberin! Cybelle!", rief Mondrianne, "Ich weiß, wer du bist, und wie leer deine Worte sind! Denk nicht, mir wäre entfallen, wie du damals meinen Lunor und sein Töchterlein bedroht hattest. Eine Mutter vergisst nicht, und eine Großmutter ebenso wenig. Wer weiß, ob er ohne deinen Überfall überhaupt am Hirschhuf gekentert wäre. Du bist bloß eine goldgierige Piratin, und nichts als Übel säst du in diese Welt."

Cybelle verdrehte ihre Augen. Mondriannes Enkelin war ihr damals etwas nahe gekommen, als Cybelle das Schiff von Mondriannes Sohn um eine Spende gebeten und ihren Säbel dabei theatralisch herumgefuchtelt hatte. Es hatte ihr nicht geschadet. Ohnehin hatte diese hammerschwingende Silberzwergin ihrem damaligen Überfall rasch ein ungemütliches Ende bereitet. Dass Kapitän Lunor kurz darauf sein Schiff in eine eigentlich gut umfahrbare Klippe gesteuert hatte, war erst recht nicht ihre Schuld. Wie nachtragend musste Mondrianne sein? Cybelles eigene Mutter hätte es kaum mit einem tadelnden Blick quittiert, wenn man ihre Tochter mit einem Schwert verletzt hätte.

Cybelle knurrte: "Wir sind keine goldgierigen \gend{Pirat*innen}{Piraten}, wir sind \gend{Freibeuter*innen}{Freibeuter}! Die größten \gend{Freibeuter*innen}{Freibeuter} der Nordmeere!" 

Unter ihrem Atem flüsterte Mondrianne: "\gend{Seeräuber*innen}{Seeräuber}!"

"Siehst du meine eleganten Kleider nicht, die gute Frisur? Mein Mantel wird von Knöpfen zusammengehalten, welche vom legendären Seekrieger Ruuf höchstpersönlich stammen, dem Beschützer der Armen und Bedürftigen. In der guten alten Zeit vor Varatans Schreckensherrschaft, als die \gend{Seekrieger*innen}{Seekrieger} sich nicht bloß um den Schutz der Reichtümer der reichsten \gend{Werftler*innen}{Werftler} kümmerten. Bevor Ruuf von Varatan korrumpiert wurde"

"Pah! Dass ihr das Symbol der gefürchteten Schwarzen Kogge hisst, verrät deine Intentionen. Deren Besatzung war ein Pack ebenso elender \gend{Pirat*innen}{Piraten} wie die deinen. Im Kampf gegen sie verlor dein so geliebter Ruuf gar sein Leben!"

Minerma war hinter Cybelle aufgetaucht und bellte: "Die Schwarze Kogge steht für mehr, als du denkst, Händlerin. Kapitän Callem war schließlich einst selbst ein Seekrieger, ehe er sich von Varatans Kriegswahn abwandte und selbst nach der Krone der Nordmeere strebte, für eine friedlichere Zukunft unter Führung der Mächte des Meeres. Wenn Ruuf doch nur auch diese Willensstärke und Weitsicht gehabt hätte."

Mondrianne lachte freudlos auf. "Callem war ein machtgieriger Mörder und die Mächte des Meeres sind nicht für ihre friedlichen Taten bekannt. Ich selbst verlor einst fast ein Schiff in Oktohans Schlund."

Minerma lachte: "Hast du dir den Riesenkraken deswegen auf die Stirn stechen lassen? Mit deinen Tätowierungen siehst du mehr aus wie eine Piratin als ich!"

Die alte Seehändlerin war wahrlich von Kopf bis Fuß tätowiert, mit kunstvollen Bildern von Sternen, rostigen Ankern, rüstigen Papageientauchern, Riesenkraken und sogar einigen grünen \gend{Zwerg*innen}{Zwergen}runen, auch wenn jene nicht magischer Natur waren. Ein unergründliches Grinsen huschte über Mondriannes Züge, als sie antwortete: "Sein und Schein, Piratin! Ihr mögt aussehen wie Edelleute und euch für \gend{Freibeuter*innen}{Freibeuter} halten, aber ihr seid immer noch genau wie die maskentragenden \gend{Pirat*innen}{Piraten} aus dem Westen. Auch du wirst nie einen Kaperbrief haben."

Jetzt war es an Cybelle, ihre Hände zu verwerfen. "Natürlich nicht, es gibt ja auch \gend{keine Meereskönig*in}{keinen Meereskönig} mehr mit der Autorität, einen Kaperbrief auszustellen. Autorität verleihen uns die \gend{Bewohner*innen}{Bewohner} der Nebelinseln, die ruhmvolle Lieder auf uns singen, je mehr Meereskreaturen wir vertreiben und je mehr Waren wir ihnen überbringen – darunter selbstverständlich auch Rum."

Mondrianne schüttelte ihren Kopf: "\gend{Eine Seekönig*in}{Einen Seekönig} über Varatanien mag es nicht mehr geben, doch \gend{Seekrieger*innen}{Seekrieger} aus Werftheim gibt es immer noch. Die werden euch einkerkern für eure Taten, sobald sie euch endlich erwischen. Und wenn man vom Seeteufel spricht ..." Mit diesen Worten langte Mondrianne an ihren Gürtel, zog ein gewundenes silbernes Horn hervor und blies hinein. Ein blechernes Tröten hallte übers Deck.

Cybelle hätte eigentlich gerne noch länger über Piraterie und Edelhaftigkeit diskutiert, doch machte ihr Mondriannes Selbstsicherheit gehörig Sorgen. Dem Hornstoß folgte das Trommeln geübten Gleichschritts auf hölzernem Deck, als die Kajütentür aufgeschlagen wurde und vier \gend{Seekrieger*innen}{Seekrieger} in ihrem typischen roten Gewand und der hölzernen Rüstung hervorsprangen. Ein Hinterhalt! Drei der \gend{Angreifer*innen}{Angreifer} trugen Naginatas, die typischen Schwertlanzen der \gend{Seekrieger*innen}{Seekrieger}. Die vierte Seekriegerin zückte hingegen gleich zwei lange Schwerter, eines von ihrem Gürtel und eines von ihrem Rücken. Sie reckte ihr Kinn provozierend in Richtung der beiden Freibeuterinnen an Bord. Ein silberner Zahn funkelte in ihrem Mund auf.

Neben Cybelle hob Minerma nervös ihren Degen und blinzelte unsicher zu Cybelle. Solche Konfrontationen waren sie nicht gewöhnt, solchen Konfrontation waren sie nicht gewachsen. War dies der Moment, Fersengold zu geben? Cybelle schüttelte kaum merklich ihren Kopf. Die Seekriegerin mit dem Silberzahn und den zwei langen Schwertern stürmte brüllend auf sie zu. Cybelle verbesserte ihre Fußstellung und machte sich auf den Kontakt gefasst.

Minerma ergriff rücklings die Flucht. 

Cybelles Säbel sauste durch die Luft und knallte gegen die erste Klinge der Seekriegerin. Mit einem lange geübten Drehmanöver fügte sie Druck auf die Hand der Gegnerin aus, woraufhin jene ihr linkes Schwert fallen ließ. Cybelle fühlte sich großartig.

Es blieb gerade genug Zeit, das Kommende zu erahnen, ehe das zweite Schwert schwer auf Cybelles eigenen Unterarm niederfuhr und ihre Welt in glühenden Schmerz tauchte. Es klimperte, als Cybelles Säbel auf das Deck knallte. Es polterte dumpf, als etwas Zweites dem Säbel folgte. Ein gehässiges Grinsen erschien auf dem Gesicht der Seekriegerin, als sie mit ihrem einen Schwert Cybelles Wange anritzte und rief: "Ergib dich, Seeräuberin, dann lasse ich dich deine andere Hand behalten."

Die Worte wurden nicht registriert. Verdutzt starrte Cybelle auf ihren pochenden Armstumpf, dessen rötliches Blut sich mit dem rauschenden Meereswasser auf dem Deck vermischte. Wie im Traum versuchte sie, Finger zu bewegen, die nicht mehr ihre waren. Ihr Sichtfeld verschwamm. Ihr Herzschlag pochte immer schneller und lauter in ihren Ohren.

Und dann war da auf einmal eine eigenartige Klarheit. 

Wenn Cybelle jetzt nicht agierte, würde sie den Rest ihres Lebens in Gefangenschaft verbringen und in einer Zwangswerft Werftheims einhändig baufällige Schiffe zusammenschustern. Diese Genugtuung durfte sie den \gend{Seekrieger*innen}{Seekriegern} nicht geben. Eher würde sie ihr Glück mit den Mächten des Meeres versuchen.

Ungeduldig griff die Seekriegerin nach Cybelles blutendem Armstumpf und setzte an, ihn auf ihren Rücken zu drehen. Es schmerzte nicht einmal. Cybelle spuckte aus, drehte ihren Arm gegen der Daumen der Seekriegerin und glitschte aus deren eisernem Griff. Und ohne eine weitere Sekunde abzuwarten, warf Cybelle sich über die Reling, ins Reich der Mächte des Meeres. Der stürmische Ozean schnellte ihr in Windeseile entgegen und klatschte ihr eiskalt ins Gesicht. Salzwasser drang in ihre Nase. Ihr Arm, der vorher nicht einmal geschmerzt hatte, brannte auf einmal, als stünde er in Flammen. Sie schnappte auf vor Schmerz, schluckte dreckige Brühe und japste nach Luft. 

Verzweifelt strampelte sie gegen die Strudel, die sie ins nasse Reich der Nerax zu ziehen versuchten. \gend{Seekrieger*innen}{Seekrieger} konnten sich dank ihrer eleganten Holzrüstungen länger über Wasser halten, doch Cybelles eleganter Freibeuterinnenmantel sog sich in Windeseile voll und erschwerte jede Bewegung – da halfen auch Ruufs Knöpfe nichts. Ihr strubbeliges Haar versperrte ihr die Sicht, aber außer Dunkelheit gab es ohnehin kaum etwas zu erkennen. Sie blubberte Luftblasen, sah verschwommen, wie diese nach rechts drifteten, und schlug kräftig nach links aus.

Ihr Kopf durchschlug die Wasseroberfläche, ihr Körper weiterhin ein Spielball kräftiger Wellen. Gierig rang Cybelle nach Luft, paddelte einarmig, blickte um sich, versuchte, blitzende Sturmwolken vor dunklen Segeln zu unterscheiden. Da schlug schon wieder eine schwere Wasserfront auf sie nieder und zog sie in die dunkle Finsternis.

Diesmal fiel es Cybelle leichter, sich zu orientieren, doch umso schwerer, die Kraft zu finden, sich über Wasser zu halten. Da vorne wirbelten dunkle Schemen im Ozean umher – harmloses Schwemmholz oder finstere Meereskreaturen? Schweren Herzens öffnete sie die Knöpfe ihres Umhangs. Mit einem Ruck löste sich das Gewicht von ihrem Hals, als der Umhang auf Nimmerwiedersehen in die Tiefen des Ozeans gezogen wurde, Ruufs Knöpfe daran. Hustend und prustend brach Cybelle durch die Wasseroberfläche. 

Ihre Köpfe ragten kaum über die Wellenkämme heraus, doch das schimmernde Haar war unverkennbar. Da waren \gend{Nix*en}{Nixen}, zwei an der Zahl, die in etwas Entfernung mit großen Augen das Geschehen vor sich beobachteten. Es geschah sicherlich nicht jeden Tag, dass über ihrer Behausung eine Seeschlacht stattfand – falls man hier überhaupt von einer Seeschlacht sprechen konnte. 

Vor Cybelle ragte der schwere Rumpf der \textsc{Capella} in die Höhe. Dahinter ritt die \textsc{Perlenechse} auf wilden Wellenkämmen. \gend{Seekrieger*innen}{Seekrieger} hangelten sich entlang gespannter Seile auf das kleinere Kaperschiff zu. Ein Seekrieger verlor den Halt am glitschigen Seil und stürzte ins kalte Nass. Prompt tauchte die eine Nixe nach ihm, obwohl er mit seiner Holzrüstung problemlos selbst hätte schwimmen können. Cybelle rief selbst nach den \gend{Nix*en}{Nixen} und winkte mit ihrem blutigen Stumpf, woraufhin jene noch eiliger Reißaus nahmen. Ein letztes Mal klatschten ihre kräftigen Fischschwänze aufs Meer, dann war Cybelles letzte Hoffnung auf Rettung verschwunden. 

Cybelle verfluchte die parteiischen \gend{Nix*en}{Nixen}. Dann sandte sie ein Stoßgebet an Kenvilar. Möge die Tückische ihre Crew vor der Zwangsarbeit bewahren. Ihr ganzer Körper bibberte. Ob vor Kälte, Angst oder Blutverlust, vermochte sie nicht zu sagen. Doch war sie lieber hier unten  in eisigen Nass als oben an Bord, wo Waffen klirrten und die Stimme einer Seekriegerin hämisch etwas brüllte. Die tiefen Überlebenschancen über Bord Gegangener waren Cybelle wohlbekannt. Land war zu weit entfernt in alle Richtungen, ihr Versteck in Klippenwacht erst recht nicht erreichbar. Ihre letzte Hoffnung war der Echsenfinger, eine zerklüftete Klippenformation, die sich weit aus dem Meer erhob und schon manches Schiff in ein Wrack verwandelt hatte – und Cybelles \textsc{Perlenechse} ihren Namen gegeben hatte. Nördlich von hier befand der Echsenfinger sich, doch wusste Cybelle nicht einmal, in welcher Richtung Norden lag.

Ein Krampf strahlte durch ihren verletzten Arm und Cybelle sank, sank immer tiefer ins gewaltige Meer. Dunkelheit und Kälte überkamen sie wie ein schützender Mantel. Ihre Lungen wollten platzen. Instinktiv atmete sie Wasser ein. Sie öffnete ihre Arme und übergab sich den Mächten des Meeres. Noch immer verspürte sie keine Furcht, bloß Kälte und Neugier, was nun geschehen würde. 

Sie erinnerte sich daran, wie ihr Vater ihr Geschichten von Gevatter Tod erzählt hatte. Von der riesigen edlen Regenbogendrächin, die in dem Moment, in dem sie ins Reich des Todes übergehen würde, ihre Flügel über sie ausbreiten, ihren Schwanz um sie winden und sie mit einem Feuerstrahl versengen solle. Alles Humbug. Und doch hielt Cybelle ihre Augen angespannt offen, versuchend, in der Finsternis des Hadrischen Meeres irgendetwas zu erkennen. Keine \gend{Nix*en}{Nixen}, keine Nerax, kein Gevatter Tod.

Ein Ruck ging durch Cybelle. Etwas rauschte und rumpelte um sie herum, noch stärker, als ihr Körper schlotterte. Es war, als wäre Cybelle von einer wilden Strömung erfasst worden. Dann wurde es merklich heller. Die Kälte zog sich zurück und eine wohlige Wärme wogte durch Cybelles Leib. Und ihr Kopf drang durch die wilde Meeresoberfläche. Luft, frische Luft, oh wie sie sie vermisst hatte!

Ein eigenartiges Ziehen in ihrem Innern gab Ruck, und jede Menge geschlucktes Meereswasser sprudelte ohne Cybelles Zutun durch ihren Mund ins Freie. Nun konnte sie hustenfrei atmen und sich umblicken. Eine unförmige, hellblau leuchtende Wassermasse umgab ihren Unterleib, sauberer, ruhiger, wärmer als der umgebende Ozean, und hielt sie eine halbe Mannshöhe über den wogenden Wellen. 

Vor ihr zeigte sich einer ihrer schlimmsten Albträume. Ihr geliebtes Schiff, die \textsc{Perlenechse}, war auf die Seite gekippt und versank im stürmischen Wasser. Ihr Heim, mit dem Cybelle so viele gute Erinnerungen verbrachte. So viele Erinnerungen, Besitz, \gend{Kamerad*innen}{Kamerad}schaft verband sie damit. Sie heulte auf.

Ein Ruck ging durch Cybelle und die magische Welle, die sie umgab. Rasch, viel rascher als selbst Minerma rudern könnte, wurde sie über das Meer gesaust, weg vom Sturm, weg von der \textsc{Capella}, auf der ihre Crew vermutlich soeben in Ketten gelegt wurde.

War dies nicht ihr Ende? Hatte sich Arkteron, der erste und höchste Sturmgeist, ihr erbarmt? War dies nichts als ein Fiebertraum, eine wohlige Illusion, während ihre Lungen sich mit Wasser füllten? Cybelles Leben war nicht vor ihrem inneren Auge vorbeigezogen. Gab es so wenig Gutes, an das sie sich erinnern könnte?

Das Pochen in ihrem Arm ließ kaum merklich nach. Dann flackerte ihre Sicht und Cybelle ergab sich der warmen, schaukelnden Umgebung.











\newpage
\section{Kirj'vana, der Meeresgeist}

Die Welle, die Cybelle mit sich mitschleppte, kam vor einer gewaltigen Klippenformation zum Halt. Mit ihrem Verständnis dieser Gefilde erkannte Cybelle die Felsen sofort: Dies war der Hirschhuf, eine von vielen Klippen, wo Cybelles Crew erbeutete Waren und allerlei sonstiges gutes Zeug versteckt hatte. 

Damals, vor langer Zeit, als die \textsc{Perlenechse} noch nicht das Symbol der Schwarzen Kogge getragen hatte. Damals, als Praeda nur ein einsamer Fischer gewesen war, der alle möglichen Sachen aus dem Ozean zog. Als Minerma tagein, tagaus in den Werften Werftheims hatte schuften müssen. Als die junge Skyld sich an der Medizinischen Akademie Werftheims schrecklichen Untersuchungen hatte aussetzen müssen, um nicht zu verhungern. Als Corsaria noch ziellos durch die Ruinen Klippenwachts gestreift war.

Die helle, warme, womöglich, nein, wahrscheinlich magische Wassermasse, die Cybelle sicher bis hier mitbewegt hatte, entließ sie auf einmal aus ihrem schützenden Griff. Einige Augenblicke lang strampelte Cybelle gegen die Meeresfluten, ehe sie mit ihrer übrigen Hand eine Furche in der Klippe zu fassen kriegte. Schwer atmend zog sie sich auf den Hirschhuf und schlotterte auf dem kalten Felsen. Sie vermisste die Wärme dieser eigenartigen Welle bereits.

Vor ihr ragte der löchrige Rumpf eines ehemaligen Einmasters empor, dem der Hirschhuf zum Verhängnis geworden war – wie so vielen Wracks vor ihm. Das Schiff war nur noch ein halbes Skelett seiner früheren Selbst. Cybelle hätte es mit einer ganzen Crew nicht mehr fahrtüchtig machen können, und ihre Crew war wer-weiß-wo. Den vermoderten Dielen zufolge war das Unglück hier nicht gerade gestern geschehen. Überlebende Schiffbrüchige gab es folglich kaum hier. Doch vielleicht noch einen Schatz? Solange es keine hinterhältigen Nerax waren, die auf leckere \gend{Plünder*innen}{Plünderer} hofften, war Cybelle zufrieden. 

\textsc{B-a-l-e-} war in großen Buchstaben auf der Seite des Schiffs zu lesen, der Rest des Namens verwittert. Ja, Cybelle konnte lesen. Ein befreundeter Silberzwerg hatte es ihr schon früh beigebracht. Etwas Weiteres, das die ehrenvolle Freibeuterin von den maskentragenden \gend{Pirat*innen}{Piraten} aus dem Westen unterschied. Sie konnte Schatzkarten und Kaperbriefe entziffern. Auch wenn es ihr momentan an beidem mangelte. Und obwohl sie nun mit ihrer fehlenden Hand wieder direkt ins Klischee einer Piratin hineinfallen würde. 

Cybelle fluchte, als sie ihre neue Realität erneut einholte. Doch als sie an sich heruntersah, hielt sie überrascht inne. Junge rosa Haut überzog den Stumpf, der einst ihre dominante Hand gewesen war. Diese magische Welle hatte sie nicht nur an einen sicheren Ort gebracht, sondern auch den Heilungsprozess beschleunigt. 

Dem Tod durch Blutverlust und Ertrinken entgangen, hieß es nun, dem Verdursten zu entringen. Cybelle nestelte lange an ihrem Gurt, bis sie einhändig den ledrigen Trinkschlauch daraus lösen konnte. Herrje, sie würde alle Handlungen neu mit ihrer linken Hand erlernen müssen. 

Ihr Trinkschlauch war schrecklich leicht. Die letzten Wassertropfen daraus vermochten es kaum, ihren Durst zu löschen. Sie hatte ihr Schiff verloren, ihre Crew, ihre Hand, bald auch noch ihren Verstand – bei allen Kreaturen der Tiefe, falls sie diese Tortur überlebte, würde ihr dieser Tag zweifelsohne als ihr schlimmster in Erinnerung bleiben.

Nach einer Runde kreativen Fluchens stemmte sich Cybelle hoch und machte sich daran, das Wrack der \textsc{Bale...} zu untersuchen. Kein Erfolg. Zum Glück wartete kein Nerax auf sie, aber auch kein Goldschatz. Und keine Waren. Dieses Wrack war bereits gründlich geleert worden, vielleicht sogar von den ursprünglichen \gend{Besitzer*innen}{Besitzern}. Hier gab es nichts mehr zu holen außer vermodertes Holz.

Cybelle kletterte über den zerklüfteten Felsen. Ein großer Spalt zog sich durch seine Mitte und teilte ihn in zwei, was diesen Klippen den Namen Hirschhuf gegeben hatte. Und am Ende dieses Spalts wartete zwar kein Goldschatz, aber zumindest ein Versteck. Damals, aus der Zeit, als Cybelles Crew zum ersten Mal ins Meer gestochen war, noch bevor sie das Segel der Schwarzen Kogge gehisst hatten. Gut verborgen in einer hölzernen Kiste unter einem grünen Tuch lag einer ihrer Notvorräte.

Der Verschluss klemmte, doch mit einigen gut platzierten Tritten hatte sich noch manches Problem lösen lassen. Der Notvorrat erhielt Trinkwasser, getrocknete Heilkräuter, gedörrte Früchte, was auch immer das Freibeuterinnenherz in Not begehren könnte. Was hatte ihre Crew damals zusätzlich hier gelassen – Schatzkarten, Ersatzteile für kaputte Schiffe, Enterhaken? 

Nein. 

Cybelle stieß ein gehässiges Lachen aus. Das Schicksal hatte einen miesen Sinn für Humor, oder die magische Welle hatte sie intuitiv zur richtigen Klippe geführt. Andächtig hob Cybelle ein gebogenes Metallstück hervor, von dem sie sich noch genau erinnerte, wie sie es ursprünglich hier versteckt hatte. Corsaria Algenbart hatte es damals besorgt, und Cybelle hatte es ausgelacht, da sie ja gerade nicht wie \gend{Pirat*innen}{Piraten} wirken wollten.

Eine messerscharfe Hakenhand.

Mit Mühen und Ringen befestigte Cybelle das Metallstück an ihrem verheilten Armstumpf. Er schmerzte nicht einmal. Mit der Hakenhand war ihr Unterarm nun sogar etwas länger vorher. Probehalber schwang sie den Haken umher. An einen Säbel kam er nie heran, doch hatte sie ihren Säbel ohnehin in der Hitze des Gefechts auf der \textsc{Capella} liegen lassen. Erneut fluchte sie die Wolken vom Himmel herab.

Nichts als das Rauschen des Meeres und das Kreischen der Seemöwen waren zu hören.

"Was war dein Plan für mich, magische Welle?!", schrie Cybelle in den wogenden Ozean. "Soll ich hier elendig zugrunde gehen?"

Prompt erhob sich eine gewaltige Wassermasse aus dem Hadrischen Meer und haute Cybelle von den Füßen. Mit ihrer verbliebenen Hand griff Cybelle vergeblich nach schartigen Felsen, doch mit der Hakenhand konnte sie sich in einer Felsspalte verkeilen und Halt finden.

Das salzige Meereswasser zog sich zusammen und blubberte zu einer menschengroßen Säule hoch, die auf einmal sehr menschenähnlich aussah. Meterlange Haarfäden aus Wassertropfen fielen über einen durchscheinenden Körper. Der Meeresgeist aus den Tiefen des hadrischen Meeres streckte plätschernd eine Hand zum Ozean aus, und eine weitere Wassersäule ließ Cybelles roten Umhang zu ihr schweben. Er war doch nicht verloren gegangen!

Cybelle wollte schon danach greifen, als ihr der verärgerte Ausdruck auf dem Gesicht des Wasserwesens auffiel. Ein fremder Geist verband sich mit dem ihren und Bilder blubberten in ihr auf. Bilder von den Weiten des Ozeans und den Tiefen des Meeresgrunds, von dunklem Wasser und dunkleren Kreaturen. Von golden glänzenden Schiffen mit weißen Segeln. Von einem schwarzhaarigen, breitschultrigen Seekrieger mit einem Mantel voller Knöpfe. 

Erneut hob der Meeresgeist den Umhang in Cybelles Richtung. Und Worte heulten aus der blubbernden Gestalt, schwer und tief, wie das Rauschen der Wellen an den Klippen. "wOhEr hAst dU dIE???"

Jetzt schlotterte Cybelle wieder, als hätte die Kälte des Meeresgrundes ihre Knochen befallen. Sie hatte selbst noch nie einen Meeresgeist getroffen, waren jene noch seltener als Geister der Erde, der Flüsse und des Windes, da die Mächte des Meeres ihre Präsenzen nur zu gerne vernichteten. Und die wenigen Geschichten, die ihr Vater ihr vor seinem viel zu frühen Tod zugebrummt hatte, als sie bloß ein kleines Mädchen auf seinem Schoß gewesen war, das waren alles üble Sagen gewesen. Erzählungen, die zeigten, dass niemand, der sich mit Meeresgeistern einließ, gut davonkam. Ihr Zorn war schwer zu besänftigen. Und Cybelle musste ja ausgerechnet mitten auf dem Ozean auf einen treffen, dem sie nirgendwohin ausweichen konnte. Diplomatie war ihr einziger Ausweg. Leider nicht eine ihrer Stärken.

"Den Mantel?", stieß sie hervor, "Den habe ich auf Sturmtal einem Salathändler abgetauscht. Nichts weiter als ein bisschen Silber aus den Schmieden Silberhalls hatte er verlangt. Dieses Silber von der verfluchten Insel ist den Taren unglaublich viel wert."

Falls das überhaupt möglich war, richtete sich der Meeresgeist noch größer vor Cybelle auf, schüttelte ihr die Mantelspitze entgegen und brüllte: "lÜÜgnErIn! wOhEr hAst dU dIE knÖpfE?!"

"Och, die ollen Dinger? Die habe ich in einem Glücksspiel einem sensiblen Typen eingenommen, in der Taverne zur Tanzenden ..."

"wElchEr tYp?! dIE knÖpfE stAmmEn vOn rUUF! wO Ist Er?!"

Genau, diese Knöpfe stammten von Ruuf, dem legendären Seekrieger, das hatte der Typ in der Taverne behauptet. Cybelles Mutter hatte ihr in ihren seltenen nüchternen Momenten gerne zugeredet, sie wäre eine Nachkommin Ruufs. Ein schöner Gedanke, auch wenn er vermutlich nicht stimmte. Halb Werftheim wollte mit dem elenden Meereskönig Varatan verwandt sein, obwohl dessen Linie sicherlich schon lange ausgelöscht war. Cybelle sah sich lieber in Ruufs Linie.

Erneut blubberten fremde Bilder vor Cybelles innerem Auge auf: Der Meeresgeist, wie er einen modrigen hadrischen Kompass an einen schwarzhaarigen Seekrieger übergab, während der Seekrieger einen Knopf von seinem Mantel an den Meeresgeist übergab. Der Seekrieger, der von diesem Kompass aufblickte, während sich vor ihm eine gewaltige Felsformation in der Form eines versteinerten Arrogs aus einem kochenden Ozean erhob. Der Seekrieger, wie er mit einer behandschuhten Hand die Wange des Meeresgeists streichelte. Ein silbern schimmerndes, beinahe unsichtbares Band, das den Seekrieger und den Meeresgeist verband.

Und Cybelle verstand. Liebe zwischen Naturgeistern und Menschen tauchten in vielen Märchen auf und endeten meistens damit, dass der Naturgeist Jahrhunderte nach dem Tod \gend{seiner Gelieb*ten}{seines Geliebten} noch trauerte. Immerhin war dieser Meeresgeist hier nicht für seine Liebe aus seinem Element verbannt worden. Nun musste Cybelle ihn beruhigen und ihm klarmachen, dass es nicht sie gewesen war, die Ruuf umgebracht hatte – sondern Seekönig Varatan, sein einst engster Freund unter den \gend{Seekrieger*innen}{Seekriegern}.

"Behalte den Mantel! Behalte die Knöpfe!", rief Cybelle, "Ich spüre, dir bedeuten sie viel mehr als mir. Doch wisse, dass ich sie nicht von deinem Liebling stahl."

Der Meeresgeist wurde noch größer und versprühte Wassertropfen in der ganzen Umgebung, ehe er noch näher zu Cybelle schwappte. Na gut, Worte waren ohnehin nicht die beste Methode, mit Geistern zu sprechen. Sie erinnerte sich zurück an den Moment, an dem sie die Knöpfe erhalten hatte. Damals, in der Taverne zur Tanzenden Tare am Hafen von Werftheim.

Ein Astzt-Kartenspiel gegen einen bulligen, vernarbten Mann. Das Zittern seines Kinns, als er die beiden Knöpfe auf den Tisch warf, weil sein Münzsack federleicht geworden war. 
Seine raue Stimme, als Rangnir knurrte: "Die Knöpfe sind bestimmt so viel wert wie dein Silber. Diese gehörten einst dem legendären Seekrieger Ruuf höchstpersönlich. Damals, ehe Varatan ihn auf Narkon verbannte." Cybelles zwei gesunden Hände, die mit Finesse und Fingerfertigkeit eine dringend benötigte Astzt-Karte aus ihrem Ärmel schüttelte, ohne dass es ihrem Gegenüber auffiel – nein, das musste der Meeresgeist lieber nicht wissen.

Das Brodeln und Sprudeln des Wassergeists wurde stiller und er schrinkte, bis er Cybelles Kopf nur noch knapp überragte. Als er erneut sprach, war seine Stimme leise plätschernd wie Wellen an einen weißen Sandstrand.

"Narkon also. Ruuf muss sich dort befinden. Danke für die Knöpfe."

Cybelle atmete auf. "Noch so gerne. Narkon ist diese riesige Insel im Nordwesten mit dem großen Silberberg – aber ich nehme an, das hast du alles schon in meinen Gedanken gesehen."

Der Meeresgeist wich von ihr zurück, ohne sich umzudrehen. Er glitt einfach rückwärts ins Meer zurück und drohte darin zu versinken. Einen Moment lang war Cybelle glücklich darüber, stellte ein wütender Meeresgeist doch eine ziemliche Gefahr für sie dar. Dann wurde ihr wieder bewusst, in welcher Lage sie sich befand, verletzt und alleine auf einer Klippe mitten im Hadrischen Meer gestrandet. Der Meeresgeist war ihre letzte Hoffnung, rasch von hier zu entkommen.

"Warte!", rief Cybelle. "Du darfst nicht einfach so nach Narkon gehen. Die Insel ist verflucht, sie würde auch dein Grab, und das hätte Ruuf sicherlich nicht gewollt! Aber ich kann dir den Typen zeigen, der mir die Knöpfe gab!"

Der Meeresgeist hielt inne.

"Sein Name ist Rangnir. Er ist ein Schiffsbauer aus Werftheim. Er hat vielleicht mehr Andenken an Ruuf. Wenn du mich an Land bringst, kann ich ihn zu dir holen."

Der Meeresgeist legte den durchscheinenden Kopf schief. Dann schüttelte er ihn, sodass sein Haar in alle Richtungen sprudelte, verspritzte noch mehr Wasser über die ohnehin schon tropfnasse Cybelle, und blubberte: "Gefährlich. Ich rettete dich kurz vor den Mächten vom Meer, aber lang ist etwas ganz anderes. Mächte mögen nicht, wenn du ihnen Versprochene aus ihrem Griff ziehst."

"Ihnen Versprochene? Die Mächte des Meeres kümmern sich doch einen Dreck um solche wie uns. Sie ließen dich mich bis hierher bringen", meinte Cybelle. "Willst du dieses Risiko nicht nochmal in Kauf nehmen? Vielleicht steht es nicht so schlecht um sie, wie du meinst."

"Glück. Nicht riskieren."

"Tja. Dann wirst du Ruuf erst recht nie wieder sehen."

Nun wandte sich Cybelle ab und schlenderte pfeifend davon. Ein Ziel hatte sie nicht, doch das musste sie auch nicht. Es dauerte in der Tat nicht mehr lange, bis der Meeresgeist sie wieder schwappend umgab und sie knurrend bis ins Wasser zog.

"Bringen wir uns vor es", blubberte er. Da konnte Cybelle auch nur zustimmen.

"Ich nehme an, du weißt, in welcher Richtung Werftheim liegt?", fragte sie, selbst im ganzen Sturmnebel kaum Ahnung habend, wo Norden lag.

"Ungefähr. Nordwesten. Wind schlecht. Aber das kann uns egal."

Und mit Saus und Braus sausten Meeresgeist und Freibeuterin über das Hadrische Meer, in eine Richtung, die hoffentlich Nordwesten war. Cybelle hatte ihre Zweifel, dass der Meeresgeist sie so leicht nach Werftheim führen konnte, hier draußen im Hadrischen Meer konnte man sich ohne Hadrischen Kompass und Seekarte rasch verirren. Aber diese Klippe würden sie umschiffen, wenn sie sie erreichten, nicht vorher.

\trennzeichen

Erneut reiste Cybelle umgeben von einer magischen Welle über das Hadrische Meer. Dieses Mal war sie allerdings wacher und erheblich weniger blutend. Zudem wusste sie nun, was es mit der magischen Transportwelle auf sich hatte. So konnte sie sich zwar entspannen und zurücklehnen, und die Reise auf der warmen Meereswelle sogar genießen, doch wurde ihr auch rasch langweilig. Nachdem sie eine Zeit lang mit ihrer neuen Hakenhand in ihrer Frisur herumgewuschelt hatte, und ihren Mantel mehr schlecht als recht wieder verknotet hatte, versuchte sie, Konversation mit ihrem stummen Reisebegleiter aufzunehmen.

"Wie soll ich dich eigentlich nennen, Meeresgeist? Hast du einen Namen?", fragte Cybelle.

Der Meeresgeist blieb still.

"Wie nannte dich Ruuf?"

"Geht dich nicht an."

"Na, dann kann ich die Zeit auch damit nutzen, dir einen Namen zu geben. Wie klingen Eabantu? Teminuhe? Freawaru? Platsch?"

Der Kopf des Meeresgeists manifestierte sich aus der magischen Welle und schüttelte sich bei jeder Erwähnung entschiedener.

"Jörm?"

"Wie die Seeschlange? Nicht ich."

Ah, jetzt sprach er plötzlich wieder mehr.

"Naryx?"

"Die Nerax sind abtrünnig, ihr Name soll nicht als Bezeichnung genutzt werden."

"Vajar?"

Stille.

"Ich komme näher, oder? Vanar?"

"Nicht ... ganz. Kirj'vana." 

Kirj'vana richtete sich etwas größer auf und proklamierte nicht ohne einen gewissen Stolz.

"Ich bin Kirj'vana, die Hüterin der Kristallkrone der starren Knochen und Weiserin des stürmischen Stroms."

"Kirj'vana, häh? Spannende Vermischung zweier Namen. Der Geist einer in der Schlacht um Klippenwacht ertrunkenen Seekriegerin namens Kirja mit einem Naturgeist namens Vanar? Ich weiß ja nicht, wie ihr Geister funktioniert, aber die Taren aus Sturmtal halten euch für im Meer Ertrunkene, nicht? Meinst du ..."

"Schweige von Dingen, die du nichts versteh", knurrte Kirj'vana Der schützende Griff des warmen Wassers um Cybelle löst sich nur einen Lidschlag lang, doch sandte der kurze Kontakt mit dem eiskalten Meereswasser bereits einen Schock durch Cybelles Magen.

Die fing sich und grummelte: "Verzeih. Kirj'vana also, behütende Hüterin und weise Weiserin von Dingen, von denen ich nichts verstehe. Freut mich, deine Bekanntschaft zu machen."

Kirj'vana gab keine Entgegnung von sich.

Fortan blieb auch Cybelle stoisch still und versuchte, Kirj'vana nicht weiter zu reizen.

Eine Zeit lang trieben die beiden weiter durch neblige Gewässer, die Freibeuterin im Meeresgeist. Unter ihnen lag für hunderte Menschenlängen dunkles Wasser, voller Spornwale und Seeschlangen. Über ihnen waren Seemöwen und Sturmwolken an einem immer weniger stürmischen Himmel. Sturmwolken wichen leichteren Schwaden. Hier und da brach sogar ein Lichtstrahl hindurch.

Es war wunderschön.

Bis Kirj'vana urplötzlich aufgellte: "pAss AUf! UntEr Uns!"

Cybelle schreckte aus ihren Tagträumen und sah sich um. Was war die Gefahr, was wollte Kirj'vana sagen, wovor wollte sie sie warnen?

Mit einem gewaltigen Grollen brach neben ihnen eine gewaltige Klippe aus dem Wasser. Es stank nach Fäulnis. Dann öffnete sich ein pupillenloses Auge und Cybelle wurde gewahr, dass dies keine Klippe war. Das war ein Arrog. Eine gewaltige Masse aus Fleisch, Fett und stahlharten Schuppenpanzern, ein Schatten eines gewaltigen Wesens aus den Tiefen des Meeres, erhob sich weit über sie. Ein Mund, groß genug um sie mit einem Happs zu verschlingen, öffnete sich und entblößte gelbe Zähne. 

"Mächte vom Meer mögen nicht, wenn man ihnen Versprochene rettet!", rief Kirj'vana.

"Zufall!", knurrte Cybelle. 

Der Arrog katapultierte seinen Oberkörper erstaunlich weit aus dem Wasser. Sein Schatten fiel über die beiden. Dann ließ sich der Arrog auf die beiden fallen. 

Cybelle verspürte ein Ziehen in ihrem Magen, als Kirj'vanas warmes Wasser sich über sie stülpte und sie unter die Wasseroberfläche zog.  "Nicht atmen. Ruhig", blubberte der Meeresgeist ihr zu, während über ihnen mit einigem Getöse der Arrog wieder ins Wasser platschte. "Vielleicht hat er nicht wir als Ziel. Vielleicht sind wir nicht gesehen."

Cybelles Ohren dröhnten, ihre Sicht verschwamm, ihre Lunge protestierte, und doch fühlte sie sich erstaunlich ruhig. Und wagemutig. Sie öffnete ihren Mund und sprach ins Wasser, so gut sie konnte: "Arrogs so nahe an der Wasseroberfläche sind meistens im Anmarsch auf eine Nebelinsel. Bring mich auf seinen Rücken. Er kann uns nach Werftheim bringen."

"Du hast Wahnsinn", ertönte Kirj'vanas Stimme als Antwort. Doch sie folgte dem Vorschlag. Erneut kehrte sich Cybelles Magen um, als Kirj'vana sie wieder nach oben beschleunigte. Die helle Wasseroberfläche schnellte auf sie zu und sie brachen direkt hinter dem gewaltigen Arrog ins Sonnenlicht. Cybelle schnappte nach Luft.

Cybelle wünschte sich ihren Enterhaken herbei, doch musste es nun halt auch ohne funktionieren. Sie löste ein letztes Stück des reißfesten Tulgori-Seils von ihrem Gurt, knotete mit zittriger Hand und Zähnen eine behelfsmäßige Schlaufe hinein und schleuderte diese dem Arrog entgegen. Das Seil klatschte gegen den mit Algen und Seesternen verkrusteten Panzer der Meereskreatur und platschte wieder ins Wasser. Cybelle stöhnte auf, da erhob sich eine magische Welle über der Wasseroberfläche und Kirj'vana befestigte die Schlaufe an einer schief stehenden Schuppe.

Sie waren dem Arrog angehängt. Und keine Sekunde zu früh, denn soeben setzte sich das Biest wieder in Bewegung. Mit einem Ruck spannte sich das Seil und Cybelle schnellte einmal, zweimal auf den Arrog zu. Wie schon zuvor das Seil klatschte auch sie gegen die stinkende grünliche Steinhaut der Kreatur. Die Berührung entlockte dem Arrog nicht einmal ein Brummen, er schien die blinde Passagierin nicht zu spüren.

Gerade noch rechtzeitig verankerte Cybelle ihre messerscharfe Hakenhand in einer Spalte zwischen zwei Rückenschuppenplatten des Ungetüms, da wand sich der Arrog umher, nach links, nach rechts, und plantschte mit seinen großen Krallenhänden ins Wasser. Doch die zuvor so leicht erschienene Beute war verschwunden. Frustriert brummte das Ungetüm und ließ sich wieder tiefer ins Wasser senken. Dann begann es wegzupaddeln, zielstrebig in Richtung Nordwesten. In Richtung Werftheim.

\trennzeichen

Als in der Ferne die ersten Lichter Werftheims sichtbar wurden, richtete sich Cybelle auf dem Rücken des Arrogs auf und machte sich für den Absprung bereit. Mit dem schweren Mantel, den sie wieder an sich genommen hatte, und der fehlenden Hand, wäre sie für sich alleine gestellt nicht weit gekommen, doch Kirj'vana bugsierte sie sanft und verlässlich an ein felsiges Ufer ein wenig abseits vom Haupthafen Werftheims.

Von dort aus hatte Cybelle einen guten Blick darauf, wie der Arrog schnurstracks auf eines der edelsten Schiffe im Hafen zuhielt und sich mit einem Brüllen darin verbiss. Mit scharfem Blick erspähte Cybelle den Schiffsnamen \textsc{Lacerta} am Rumpf. Nun, die elegante \textsc{Lacerta} wurde rasch von Arrogklauen zerfurcht. Üblicherweise würde sie sich schlecht fühlen beim Gedanken an die schöne Arbeit der \gend{Schiffshandwerker*innen}{Schiffshandwerker}, die nun verloren ging. Aktuell fühlte sie allerdings gar nichts, außer Genugtuung. Schon ertönten die ersten Alarmglocken, und fackelschwingende \gend{Seekrieger*innen}{Seekrieger} bemannten Ballisten, um das Ungetüm von der Hauptstadt zu vertreiben. Möge es einige von ihnen mit sich reißen.

Cybelle versteckte sich hinter einigen Fässern am Pier. Land, süßes Land! Sie hatte es geschafft.

Als sie hinter sich blickte, bemerkte sie, dass Kirj'vana ihr nicht gefolgt war. Der Meeresgeist hielt sich im seichten Wasser und wagte nicht, an Land zu treten. Zur durchscheinenden Person zurückkehrend, versprach Cybelle: "Ich komme wieder. Wir werden herausfinden, was es mit diesen Knöpfen auf sich hat." Sie hatte keinerlei Intention, dieses Versprechen zu halten, sofern sie nicht persönlich etwas daraus gewinnen konnte. "Und wer weiß, vielleicht kannst du uns ja bald wieder übers Wasser transportieren."

"Nie wieder", sprach Kirj'vana entschieden. "Du hast Arrog gesehen. Mächte vom Meer sind nicht am Spaßen."

Sie strich mit einer wässrigen Hand über Cybelles Gesicht. Tropfen perlten aus ihren Haaren und ihrem Mantel und flossen zum Meeresgeist zurück, bis Cybelle beinahe trocken am Ufer stand. Eine unergründliche Traurigkeit lag in Kirj'vanas Augen.

Auf ihrem Weg zu den Häusern Werftheims blickte Cybelle hin und wieder zurück. Kirj'vana stand jedes Mal immer noch unbewegt im seichten Wasser und wogte still vor sich hin.





\newpage
\section{Rangnir, der Schiffsbauer}

Zitternd stolperte Cybelle durch das nächtliche Werftheim. Sie war trocken, doch änderte dies nichts an der beißenden Kälte. Ihr Magen knurrte. So menschliche, so alltägliche Bedürfnisse verblassten gegenüber den wilden Tagen, die sie hinter sich hatte. Ein Schiff, eine Crew, eine Hand hatte sie verloren. Ein Meeresgeist aus den Tiefen des Ozeans hatte sie gerettet und wartete auf einen Bericht über ihren verfluchten Liebhaber. Einen Arrog hatte sie geritten. Und nun hatte sie Hunger.

Der Duft von frischen Drachenbohnen zog sie zur Taverne zur Tanzenden Tare am Haupthafen Werftheims. Versteckt lungerte Cybelle in einer Seitengasse und zuckte bei jedem zu laut bellenden vorbei streunenden Hund zusammen. Der Anblick der berühmten Werftheimer Fischpfanne ließ ihren knurrenden Magen noch mehr schmerzen. Jenes Gericht hatte es sogar in die andorischen Kochbücher geschafft.

Die Taverne zur Tanzenden Tare war nicht das bestbesuchte Gasthaus von Werftheim. Dies wäre "Varatans Thron" in der Oberstadt. Doch hatte jener Fakt mehr mit der relativ kleinen Größe der Taverne zu tun statt mit mangelnder Beliebtheit. "Varatans Thron" stand großflächig auf einem prachtvollen Grundstück, auf dem einst Seekönig Varatan selbst gelebt hatte. Die Taverne zur Tanzenden Tare stand hingegen eingequetscht zwischen anderen Hafenhäusern am Wasser, und die engen Gassen füllten sich jeden Abend mit Speisenden, die im Haus selbst keinen Platz mehr gefunden hatten.

Die Taverne hatte eine berühmte Geschichte. Sie handelte vor langer Zeit, als Werftheim noch Holtpoort genannt worden war. Denn dies war selbst lange gewesen, bevor Seekönig Varatan die Insel mit Kriegsschiffen aus Klippenwacht erobern und ganz selbstgerecht in "Varatanien" umbenennen würde. Angeblich hatte es zu jener Zeit eine rebellische Tare namens Zengin gegeben, die ihr heimatliches Sturmtal hinter sich gelassen hatte, obwohl Taren Sturmtal typischerweise nie verließen. 

Die großen Tanzfeste von Sturmtal waren heutzutage wohlbekannt und gut besucht. Zu Zengins Zeiten waren diese jedoch nur Männern vorbehalten gewesen. Frauen und andere Taren waren nicht zugelassen worden, da dies den machtvollen alten Männern obszön erschienen war. Und so hatte die tapfere Tare Zengin die gewagte Reise auf einem kleinen Boot nach Holtpoort auf sich genommen, um auf einer freieren Insel ihr Glück zu suchen.

Holtpoort war schon damals ein großer Warenumschlagsplatz gewesen. Zengin hatte ausgelassen getanzt zu Musik der wildesten \gend{Bard*innen}{Barden}, zu welcher sie zuvor nur im Schatten ihres dunklen Zimmers hatte tanzen dürfen. Zengins fünf Mägen hatten sich mit viel zu viel Rum und Schnaps gefüllt, denn der tanzenden Tare waren diese Mittel gänzlich ungekannt gewesen – damals war Rum in Sturmtal noch weniger verbreitet als heute, hatten die Taren doch besser als die umliegenden Inseln das darin enthaltene Gift als solches erkannt. Taren mochten gutes Quellwasser ohnehin lieber, das gab ihnen Stärke. 

Jedenfalls kam es so, dass eine sturzbesoffene Zengin mit der Schnapshändlerin Stybak zusammenstieß und ihren ganzen Karren umriss. Mehrere Flaschen besten Rums gingen dabei zu Bruch. Und Stybak verliebte sich augenblicklich in Zengin. Zengin hingegen nahm Stybak zunächst kaum wahr, bis die Schnapshändlerin selbst zu tanzen und singen begann. Dies nicht zwingend gut, aber tempramentvoll und unbeschwert. Viel Rum füllte die Mägen, ehe der nächste Morgen anbrach. Zengin versprach bei ihrem Abschied, bald wiederzukommen. 

Und das tat sie auch, nüchterner und zurechnungsfähiger, und die beiden verbrachten viele frohe Abende. So beschloss die Schnapshändlerin Stybak, hier in diesem Viertel zu bleiben, einen kleinen Stand aufzutun, und wilde Tanzgelage an der Seite ihrer geliebten Tare Zengin zu führen.

Und diese waren so beliebt, dass der Stand bald zu einer kleinen Hütte wurde. Viele Reisende kamen vorbei und erzählten, was hier und dort auf den tausend Inseln des hadrischen Meeres geschah. Und so wurde bald aus den Überresten des alten Karrens die Taverne zur Tanzenden Tare. Noch heute wurden in großen Saunen hinter der Taverne traditionelle Tänze aus dem Sturmtal für Interessierte aller Geschlechter aufgeführt. Und noch heute wurde der Schnaps "Tarenblut" in der Schenke in Gedenken an Zengin serviert.

\trennzeichen

Neben dem Taverneneingang hing ein Schild mit verschiedenen Aushängen zu Arbeitsangeboten, Marktstandpreisungen und gesuchten Personen.

Ein Anhang fiel Cybelle sofort ins Auge. Ein Stück Pergament, auf dem ihr Schiff mehr schlecht als recht abgebildet war, neben einer Skizze von Cybelle höchstpersönlich. In jener hielt sie einen Anker in der Hand und gute Güte, hatte die Person, die dies gezeichnet hatte, je einen echten Anker gesehen? Das war ein Stockanker, der Ankerstock durfte doch nicht parallel zu den Fluken platziert werden, so würde er doch gar nicht funktionieren! Einen Bart hatte ihr jemand auch angemalt, doch fand Cybelle dies weniger tragisch. Sie schlug ihre Kapuze höher und wartete.

Langsam, langsam verließen die letzten Gäste die Taverne und drinnen verlöschten immer mehr Kerzen. Cybelle fragte sich bereits, ob sie sich geirrt hatte und ihr Ziel den Abend woanders verbracht hatte.

Da vernahm sie Rumoren aus dem Inneren: "Nein, du hattest genug Rum! Mit solchem Gebaren vergraulst du mir die Kundschaft für Tage!"

Die Türe öffnete sich und ein Trunkenbold wurde herausgestolpert.

Ein kleiner, breit gebauter Mann, der mit einer leeren Flasche herumfuchtelte und vor sich hin brummelte. Ihn hatte Cybelle schon so oft gesehen, dass es sie schmerzte, ihn in diesem Zustand zu erblicken.

Nun trat ein bierbäuchiger blauhäutiger Bursche mit buschigem Bart in die Nacht. Breat, der wohlbekannte Wirt dieses wohligen Wirtshauses höchstpersönlich. Er schleppte den Trunkenbold zur nächsten Tränke und pflatschte ihm kaltes Wasser ins Gesicht. 

"Kehre nicht wieder zurück, bis du dich besser im Griff hast, hä? Und wasch dich, du stinkst nach Einhorndung. Ich bin da Experte."

Breat sich die Hände an seiner Schürze trocken und kehrte zu seinem Gasthaus zurück. Der Trunkenbold zog sich indes an der Tränke hoch. Er schaffte es einige Schritte die Straße entlang, ehe er in eine Pfütze stürzte und sich vollschlammte. Er murmelte zu niemand bestimmten etwas von der heilenden Wirkung einer Einhorndungkur.

Cybelle schlenderte auf ihn zu. Unauffällig fischte sie einen versilberten Kurzdolch aus bestem Zwergenstahl aus ihrem durchnässten Mantel.

"Na, wen haben wir denn hier?", rief sie laut.

Der Trunkenbold zuckte zusammen, als ihr Dolch auf ihn niederfuhr. Dann erklang ein helles Ringen, als Cybelles Dolch auf seinen traf. Denn er hatte ebenso rasch einen identischen versilberten Kurzdolch aus bestem Zwergenstahl gezückt und ihr entgegengehalten. Ein leises Summen ging durch die beiden Dolche, die vor sich hin vibrierten und sich weigerten, einander zu berühren.

Mit aufgerissenen Augen murmelte der Trunkenbold heiser: "Die Dolche der Verbundenheit verbinden. Cybelle, du bist es wahrhaftig. Du bist zurückgekehrt."

"Habe ich es dir nicht versprochen, Rangnir?", meinte Cybelle. "Eines Tages komme ich zurück. Eines Tages nehme ich dich mit auf ein Abenteuer. Eines Tages lasse ich dich deine langweiligen Schiffsbautage zurücklassen. Dieser Tag ist heute."

Rangnir ließ sich hochziehen und von Cybelle zu seinem Haus führen. Wobei Cybelle dessen Position aus dem immer mehr vor sich hinbrabbelnden Rangnir fischen musste und "Haus" zu viel versprach. Es war mehr ein feuchter Keller unter einem baufälligen Gebäude, in dem noch mehr \gend{Werftheimer*innen}{Werftheimer} hausten, die ihren Lohn bei Glücksspiel und Trank verspielt hatten.

Rangnir ließ sich direkt auf einen weichen Sack fallen und schnarchte los, die leere Flasche noch immer fest in seiner Hand. Cybelle kannte sich an diesem Ort nicht aus, doch fand in kurzer Zeit einen Sack voller Honignüsse und getrockneten Tangs, an dem sie sich gütlich tat.

Der dringendste Hunger gestillt, fand sie unter einem Tuch einen Holzstapel, von dem sie ein bisschen Holz zur Feuerstelle brachte. Das Feuer zu entzünden, kostete noch weitere Zeit, während sie sich mit Feuerstahl und ihrer schmerzenden Hakenhand abmühte. Aber auf einmal funktionierte es. Daran konnte sie sich gut aufwärmen.

Dann verscheuchte sie einige Hühner von einem Strohhaufen, ließ sich darauf nieder und wanderte ins Reich der Träume. Sie vermisste ihre Hängematte in der \textsc{Perlenechse}. Und sie fürchtete die kommenden Albträume.

\trennzeichen

Cybelle träumte nicht, oder zumindest keine Träume, an die sie sich erinnerte. 

Es roch lecker, als sie aufwachte.

Sonnenstrahlen streichelten ihr Gesicht durch die kleinen vergitterten Kellerfenster. Jemand hatte eine schwere, trockene Decke über sie gezogen, die ihr etwas Wärme und Geborgenheit schenkte. Sie bibberte nicht, überraschenderweise. Das Feuer in der Feuerstelle brannte noch, was sie ebenfalls überraschte. Jemand hatte es gefüttert.

Rangnir war bereits wach, in sauberere Kleidung gekleidet, und kratzte in einem Topf über dem Feuer an etwas Essen herum.

"Mist, du bist bereits erwacht!", knurrte er, "Ich wollte dich eigentlich überraschen. Na, die Eier sind dennoch gut. Bleib einfach liegen. Erhole dich. Es sieht so aus, als ob du es nötig hättest."

Er war wie ausgewechselt. Fort war der betrunkene Tölpel von gestern. Das hier wirkte schon mehr nach dem Rangnir, den Cybelle kannte. Im Nu hatte er Cybelle ein üppiges Frühstück bereitet, mit frischem Brot, Rührei, geräuchertem Seetang und Werftheims bestem Kräutertee.

Auf einmal fühlte sich das feuchte Kellergewölbe ein bisschen mehr nach einer Heimat an.

"Ich muss mich entschuldigen für gestern", rumpelte Rangnir, während er selbst an seinem Frühstück kaute, "Ich zeige nicht meine beste Seite, wenn ich zu tief ins Glas guckte. Passiert mir leider allzu oft. Ich muss daran arbeiten. Danke, dass du mich hierher zurückgebracht hast."

Cybelle antwortete nicht. Sie hatte so etwas schon zu oft von ihrer Mutter gehört, als dass sie eine echte Verhaltensänderung erwarten würde.

Rangnir blinzelte zu ihrer Hakenhand, dann zurück auf den Tisch.

"Ich habe gehört, dass die \gend{Seekrieger*innen}{Seekrieger} Corsaria gestern in die Zwangswerft eingewiesen haben. Den Rest eurer Crew vermutlich auch bald. Es sieht nicht gut aus um sie."

Cybelle schwieg missmutig. Das hatte sie befürchtet.

Rangnir sah sie nachdenklich an. "Cyb, was ist geschehen?"

Und Cybelle erzählte es ihm. Alles, was schiefgelaufen war in den letzten Tagen. Sie brachte die Fakten emotionslos zu Tage, redete rasch, um sich nicht von aufwallenden Gefühlen überwältigen zu lassen, und ließ dennoch keine relevanten Details aus. Rangnir hörte ihr aufmerksam zu, stellte hin und wieder Zwischenfragen, und schüttelte immer heftiger seinen Kopf. 

Dann war es vorbei. Cybelle nahm einen tiefen Schluck ihres Kräutertees. Ihr Kinn zitterte. Rangnir setzte an, ihr beruhigend auf die Schulter zu klopfen. Er ließ die Hand allerdings wieder fallen, als er Cybelles finsteren Blick sah.

Lange blieb es still, bis Rangnir brummelte: "Ich nehme nicht an, dass du die Kohle hast, um dir einen Besuch bei den hochnäsigen Physikussen der Medak zu leisten. Aber den Sprung nach Sturmtal könntest du überstehen. Ein bisschen Aufenthalt im Sturmtal hat noch jedes vergängliche Leiden kuriert."

"Sonderlich vergänglich ist eine fehlende Hand nicht", meinte Cybelle. "Wenn ich bloß diese Seekriegerin in die verbliebenen Finger kriege ..."


"Dann vielleicht ein Aufenthalt im Untergrund? Du solltest dich verdeckt halten, doch ich könnte zur Halle der Hafenmeister gehen in deinem Auftrag. Mir wurde gesagt, Sainus Teren Maltus habe ein offenes Ohr und viel Einfluss. Die Taverne zur Tanzenden Tare ist bloß seinetwegen fast so gut beliebt wie 'Varatans Thron' in der Oberstadt – und das, obwohl der Wirt Breat ein Arranier ist, einer von Kenvilars Blute. Er lässt sich sogar manchmal von einem jungen Verfluchten aushelfen, und die Kundschaft hat sich daran gewöhnt. Da wärst du auch mit einer Hakenhand nicht fehl am Platze."

"Denkst du, ich will in einer Spelunke Kupferstücke verdienen, darauf wartend, dass mich jemand erkennt?", spuckte Cybelle aus. "Während meine Crew gezwungen wird, weitere Kriegsschiffe herzustellen?!"

Erneut schüttelte Rangnir seinen Kopf.

"Was ich brauche, Rangnir ... was wir brauchen ... ist ein Schiff."

"Ich nehme nicht an, dass du viel Gold hast?"

Cybelle schüttelte ihren Kopf.

Rangnir brummelte vor sich hin. "Das heißt Stehlen. Das heißt, meinen Zugang zu den Schiffswerften ausnutzen. Das heißt, meine Stellung als ehrenwerter Schiffsbauer verlieren."

"Und wie glücklich bist du als ehrenwerter Schiffsbauer? Wie gut kommst du über die Runden? Hast du Kameradschaft? Trinkst du deswegen immer noch Nacht für Nacht alleine in einer Spelunke?"

Rangnir schwieg. Dann griff er nach seinem Mantel. "Komm. Gehen wir spazieren."

Die beiden schlenderten zurück zum Hafenviertel. Rangnir lief stolz herum, wurde links und rechts von Leuten begrüßt oder zumindest zugenickt. Cybelle hatte ihren Mund durch ein Tuch verdeckt, wie es die Fischer auf ihren windigen Abenteuern in der kalten See trugen, und lief krumm, um älter zu wirken.

Der Hafen war gut besucht. Werftheim war schließlich der größte Warenumschlagplatz des Nordens. Verschiedenste Schiffe waren hier angelegt oder am Ein- und Ausfahren. Menschen wuselten herum, Feuer in verschiedenen Farben wiesen zur kollisionsfreien Ein- und Ausfahrt aus. 

Das Schiff, das der Arrog bei Cybelles Ankunft umgerissen hatte, war etwas abseits an Land gezogen worden. Schwere Furchen von Arrogklauen zierten seine Seite. Verschiedene \gend{Baumeister*innen}{Baumeister} inspizierten die Schäden. Trotz letzteren war das Schiff klar als zugleich wunderschön und windschnell erkennbar. \textsc{Lacerta} stand in geschwungenen Lettern am Rumpf. Cybelles Hunger war geweckt.

"Wie lange, bis es wieder fahrtüchtig ist?", fragte Cybelle.

Rangnir starrte sie ungläubig an. "Das hier?!"

Cybelle nickte begeistert: "Sogar mit einem Sturmsegel! Damit sind wir dem Willen von Arkteron, dem Herrn der Stürme, nicht mehr hilflos ausgeliefert!"

Rangnir guckte hingegen sehr unsicher drein. "Das ist die \textsc{Lacerta}! Dieses Schiff gehört Seemann Emmett höchstpersönlich. Seine Wut wäre noch schlimmer als die eines armen Schluckers, da er die Mittel hat, Diebe lange verfolgen zu lassen."

"Umso besser! Weil Emmett viele Schiffe hat, wird er eines vielleicht nicht speziell vermissen! Insbesondere eines, das ein Arrog bereits angefressen hat."

Rangnir schluckte schwer: "Ich kann dafür sorgen, dass ich beim Reparaturprozess dabei bin. Und wenn es fertig ist, kann ich mich bei dir melden. Aber Cyb, bedienen kann ich ein Schiff nicht. Bauen und Nutzen sind zwei sehr verschiedene Dinge. Und drei linke Hände und eine Hakenhand reichen ohnehin nicht. Wir brauchen mehr Leute."

"Lass das nur meine Sorge sein. Eine gute Kapitänin lässt ihre Crew nie im Stich. Und das wird unsere nächste Station."

Cybelle ließ sich zum Kerkerturm führen. Ein schmuckloses, hässliches Gebäude. Die Holzlatten, die es zusammenhielten, waren mit Eisenstäben verstärkt.

Ein lautes Horn ertönte. Eine Reihe von \gend{Insass*innen}{Insassen} verließen das Gefängnis in Reih und Glied, ihre Arme und Beine gefesselt mit sehr dünnen Seilen, die ihnen zweifellos ins Fleisch schnitten.

Da vorne lief Corsaria Algenbart, den Kopf stolz erhoben, eine neue Narne an ihrer Wange. Da hinten der kleine Praeda, die Schultern niedergeschlagen gesenkt. Minerma, die direkt hinter ihm stolperte und ruppig von einer Seekriegerin wieder in die Höhe gezogen wurde. So viele ihrer Crew erkannte Cybelle. Aber auch so viele andere \gend{Insass*innen}{Insassen}. Und sie wurden von \gend{Seekrieger*innen}{Seekriegern} in allen Richtungen umrundet. 

Sie alle watschelten zur Schiffswerft, wo sie beim Schiffsbau oder auf den Tangfarmen zwangsarbeiten würden. Keine unmittelbare Gefahr, keine fehlenden Finger, keine Anzeichen von Auspeitschungen. Für den Moment ging es ihnen den Umständen entsprechend gut. Sie konnten einige Wochen in diesem Zustand überleben. Das gab Cybelle und Rangnir genügend Zeit, einen Fluchtplan zu überlegen.

Als ihr Blick zum Meer schweifte, glaubte sie, für einen Moment eine menschliche Gestalt über der Gischt schweben zu sehen. Kirj'vana hatte sie nicht vergessen. Aber sie musste jetzt warten.

Cybelle Blick schweifte zum Kerkerturm zurück. Stabile Gitter machten alle Fenster undurchdringlich, wenn man nicht gerade ein Wassergeist war. Doch auch ein Wassergeist konnte niemand Menschliches aus dem Kerkerturm retten. Gewisse Artefakte allerdings schon.

"Wohin nun?", fragte Rangnir. 

"Zurück zur Taverne. Ich muss meinen Mantel zurückholen."

"Hast du Breat einen geliehen? Du hast doch schon einen schönen."

"Durchaus. Aber es geht nicht um einen gewöhnlichen Mantel. Sondern einen wunderbaren."

\trennzeichen

Im Tageslicht sah die Taverne zur Tanzenden Tare nicht mehr so heruntergekommen aus.

Cybelle hatte die Kapuze hochgezogen und drückte sich in die Schatten der Seitengasse. Sie notierte sich alle möglichen Fluchtwege, sondierte alle vorbeischlendernden Personen nach Gefahrenpotential. Niemand achtete auf sie und Rangnir.

Stundenlang warteten sie dort und beobachteten ein- und austretende Gäste der Taverne. Der Wirt ließ sich Zeit.

Doch dann, endlich, als die Sonne schon tief am Himmel hing, schwang Tür zur Taverne zur Tanzenden Tare schwang auf, und der Wirt Breat trat heraus. Auf seinen Schultern trug er einen kleinen Wichtel, aus dessen langen Ohren weiße Haarbüschel wuchsen, und der soeben dem Wirt zukrächzte: "Wirklich, mit deinem Notfall-Mera-Stein ist alles in Ordnung. Gelegentliches Abstauben würde ihm nicht schaden, aber das Ding läuft, wie es soll. Ich suche dich sonst in einem Jahrhundert wieder auf, um ihn erneut zu testen."

"Mein lieber Wrort, bis dahin bin ich schon längst im Reich der Toten", meinte Breat seufzend. "Aber die Taverne zur Tanzenden Tare wird weiter bestehen, so die Mächte des Meeres ihr wohlgesinnt sind, von dem her wird \gend{meine Nachfolger*in}{mein Nachfolger} dich sicher gerne empfangen. SANTA!"

Auf seinen Ruf rannte ein bleicher, schlaksiger Junge von vielleicht neun oder zehn Sommern herbei. Seine Haare leuchteten in einem grellen Weiß und seine Augen waren hinter dunkelgrünen Augengläsern verborgen.

Breat setzte den Wichtel ab. "Bring ihn zur Medak", knurrte er dem kleinen Jungen zu. Er schnippte ihm eine Kupfermünze zu.

Der Junge brummelte unter seinem Atem etwas von Gier und Geiz grandiosen Ausmaßes, woraufhin ihm der Wichtel eine größere Münze zusteckte. Breit grinsend sprach der Junge lauter zum Wichtel: "Zur Medizinischen Akademie derer zu Holtpoort, also? Wrort war der Name? Wollen die \gend{Physik*as}{Physikusse} Euch untersuchen? Ich bin übrigens Santalion."

Während Santalion mit Wrort im Schlepptau davoneilte, krächzte der Wichtel ebenso atemlos: "Freut mich überaus, ja, ich bin Wrort. Die \gend{Physik*as}{Physikusse} würden meine Spezies sicher gerne untersuchen, aber vielmehr will ich für sie Licht ins Dunkel bringen und ihnen ein paar Tatsachen aufdecken. Eigentlich führte meine Reise mich aus ganz anderen Gründen hier nach Werftheim, aber nach dem, was Breat mir zu ihren Aderlass-Praktiken erzählte, kann ich hier sehr effektiv eingreifen. Dafür bräuchte ihn jedoch noch einige Instrumente. Ich weiß, die Chance ist klein, aber könnte man hier in der Nähe womöglich ein technomagisches Gelumeter auftreiben?"

Santalions ungläubige Antwort kam wie aus der werftheimschen Arcuballiste geschossen: "Das hier ist Werftheim! Für genug Geld gibt es hier alles." Und schon waren die beiden um die Ecke verschwunden.

Der Wirt Breat zog sich indes in seine Taverne zurück. 

"Gehe jetzt, wo Breat noch so nahe ist!", zischte Cybelle, und stieß Rangnir nach vorne. 

Rangnir stolperte über die Türschwelle der Taverne und Breats brummige Stimme wurde laut: 
"Ahhh, Rangnir, ich dachte schon, ich hätte dich gestern völlig vergrault. Das übliche für dich? Komm, ich war gestern etwas grob, aber Eierschwamm darüber ... und wenn wir schon von diesen sprechen, hätte ich heute eine feinste sämige Sauce mit ..."

"Heute nicht", sagte Rangnir. "Ich bin hier wegen Cybelles Mantel."










\newpage
\section{Alfreck, der Ausbrecher}

Es war ein handelsüblicher Tag am Kerkerturm gewesen.

Die Unfreien hatten lange und hart gearbeitet, in den Werften und an den Tangfarmen. In Reih und Glied waren sie zurück in ihr Gefängnis befördert worden. Nun war es Abend. Nach dem gewöhnungsbedürftigen Futter wurden die Spielkarten ausgepackt.

Alfreck las seine Gegenüber wie offene Bücher. Er war sehr gut darin, andere Menschen zu lesen, das war schon immer seine Stärke gewesen. Natürlich gewann er nicht jede Spielrunde. Er achtete darauf, seine \gend{Mitspieler*innen}{Mitspieler} bei der Stange zu halten, ihnen oft kleine Gewinne und hin und wieder sogar einen Hauptgewinn zu lassen, damit sie weiter mit ihm spielten. Über kurz oder lange würde natürlich er der Sieger sein. Dies war das wahre Spiel für ihn. Ein Zeitvertrieb, der die klein wenige Macht, die er in der Zwangswerft hatte, mehrte. Und der ihm unglaublich Spaß bereitete.

Selbstverständlich lamentierte er währenddessen immer wieder in seiner hohen Stimme vor sich hin. Dass er rauswollte aus dem Kerkerturm. Dass seine arme Tochter da draußen in einem Waisenheim auf ihn wartete, seit ihre Mutter gestorben war. Dass er doch nur hier raus und sich um sie kümmern wollte. So ganz wehleidig, dass nur die Hälfte der \gend{Insass*innen}{Insassen} auf die Idee kam, dass er sie anlügen könnte, und dass noch wenigeren auffiel, wie sehr er sie währenddessen im Kartenspiel ausnahm. Genug, um Mitleid oder gar Häme zu wecken, und doch differenziert genug, dass er seine Gegenüber nicht vergraulte.

Er langte sich an den kahlen Kopf, rieb sich die grauen Augen, verwarf seine Hände und warf seine Karten von sich. Leider wurde seine Darbietung heute nicht genug gewürdigt. Der kleine Praeda war angespannter als sonst. Er achtete noch schlechter auf seine Spielkarten, verlor seine letzten Socken an Alfreck und regte sich nicht einmal darüber auf. 

Alfreck wäre typischerweise davon ausgegangen, dass Praeda eine Magenverstimmung hatte. Das kam nicht selten vor bei der Tage alten wässrigen Brühe, die hier als "Suppe" ausgeschenkt wurde. Aber es doch steckte mehr dahinter. Praeda guckte immer wieder verstohlen gen Himmel. Und nicht nur er. Corsaria Algenbart hatte sich, obwohl sie um diese Zeit üblicherweise längst in ihrer Zelle war und sich der Pflege ihrer tangigen Haare widmete, auf einer Sitzbank niedergelegt und betont die Augen geschlossen – aber hin und wieder flatterten ihre Lider. Und zwei weitere Leute fielen Alfreck auf. Interessanterweise allesamt aus der Gruppe von Neuankömmlingen, die vor einigen Wochen eingeführt worden waren. \gend{Pirat*innen}{Piraten}, die bei einem Überfall auf das Handelsschiff \textsc{Capella} der Seehändlerin Mondrianne eingefangen worden waren. Sie alle hielten dem Himmel im Blick.

Alfreck notierte es sich und schielte ebenfalls hin und wieder in die Wolken.

Eine große Möwe flatterte über ihnen und krächzte dreimal. Einige Herzschläge lang flatterte sie still vor sich hin, dann kamen vor ihr nochmals drei Krächzer.

Praeda ließ seine Astzt-Karten fallen. Er murmelte etwas davon, dringend auf die Toilette zu müssen, lief dann jedoch in die entgegengesetzte Richtung davon und schubste Corsaria an.

Diese wartete genau eine Minute, ehe sie sich aufschwang und zu den Wachen glitt.

"Lia! Deren! Stimmen die Gerüchte von einem neuen Würfelspiel aus Hadria, das sich unter den Wachen verbreitet? " Und Corsaria fing ein Gespräch an mit den beiden \gend{Seekrieger*innen}{Seekriegern}, die den Hof im Blick behalten sollten. Lias kehliges Lachen war über den ganzen Hof zu hören.

\gend{Idiot*innen}{Idioten}! Lias und Derens Schicht würde gerade zu Ende sein, und die nächste Schicht tauchte manchmal früher auf, um ihnen ein etwas weniger spätes Abendbrot zu ermöglichen. Unberechenbar. Wenn sie eine Ablenkung wollten, war jetzt der schlechteste Moment, man müsste doch bis zur nächsten Schicht abwarten!

Doch was auch immer Praeda, Corsaria und die anderen Neuankömmlinge geplant hatten, wartete nicht. Die große Möwe kam erneut angeflattert. Diesmal tiefer, langsamer, mit etwas unförmigen Großem in ihren Klauen. 

Der kleine Praeda starrte sie direkt an. Corsaria hatte noch immer die zwei Wachen in ein Gespräch verwickelt, aber da tauchte deren Ablösung auch schon auf.

Alfreck beäugte heimlich das Paket, das die Möwe brachte, versuchte abzuschätzen, wie wertvoll es sein könnte.

Dann beschloss er, es zu riskieren.

Unter dem Tisch schüttelte Alfreck eine gezinkte Karte aus seinem Ärmel und schleuderte sie mit einem gezielten Kartenwurf unter Kargos Stuhl. Dann ächzte er sich mit einem "Habt Dank für die milde Gabe, meine Spielgenossen. Ach, wie sehr ich mir wünschte, statt mit euch elendem Pack mit meiner Tochter zu spielen" hoch und setze an, sich vom Tisch zurückzuziehen. 

Die einäugige Storve protestierte: "Komm schon, nur eine Runde mehr! Praedas Einsatz könnte deiner werden." 

Alfreck winkte ab und wanderte um den Tisch davon. Aus sicherer Entfernung deutete er unter Kargo und sagte unschuldig: "Sammelt eure Karten ein, unser einziges gutes Kartendeck muss vollständig bleiben."

"Es ist vollständig!", grummelte Kargo.

Da hatte Storve bereits unter Kargos Stuhl gelangt und Alfrecks gezinkte Karte hervor gefischt. Das goldene Astzt, die wertvollste Karte im Spiel. Einzigartig, obwohl alle wussten, dass Storve gerade erst das goldene Astzt gespielt und in den Kartenstapel zurück gemischt hatte. 

Storve war nicht gerade für ihre feine Art bekannt und brüllte Kargo als Falschspieler an. Jener fackelte nicht lange und nahm Storve in den Schwitzkasten.
 
Alfreck spazierte davon, weit weg vom Getümmel und Geschrei, und die beiden neuen Wachen rannten geradewegs an ihm vorbei auf den Klamauk zu, statt sich um die Möwe zu kümmern. Die Möwe ließ indes ihr Paket fallen. Praeda lief unschuldig darauf zu, doch Alfreck war schneller.

"Sieh an, was haben wir denn hier?", grinste er.

"Gib das her!", zischte Praeda.

Alfreck befühlte das Paket. Es war weich. Keine Waffen, wie er gehofft hatte. Aber auch keine Nahrung. Ein Kleidungsstück? Es fühlte sich merkwürdig widerständig an.

"Was ist das?", fragte er Praeda unschuldig.

Praeda druckste etwas herum, dann sagte er: "Unser Weg nach draußen."

Alfreck antwortete: "Lasst mich Teil davon sein."

\trennzeichen

Das Hadrische Meer plätscherte sanft vor sich hin.

Im Schatten der Werft ankerte die reparierte \textsc{Lacerta}. Ein stolzes Schiff. Das Sturmsegel glänzte leicht im Sternenlicht, doch der Mond blieb aus. Perfekte Konditionen, um das Schiff zu entführen, ehe es morgen offiziell wieder Teil von Seemann Emmetts Flotte geworden wäre. 

Ihr Herz schlug bis zum Hals. Rangnir hingegen schien ganz ruhig. Langsam und stetig schlug er mit seinen Paddeln ins Wasser und brachte das kleine Ruderboot näher an die \textsc{Lacerta}. 

Cybelle kletterte an Bord. Wachen schien es keine zu geben. Sie pfiff nach Rangnir, welcher ihr ein unförmiges Paket hochreichte. Mit ihrer Hakenhand löste Cybelle den Knoten. Blauer Stoff kam zum Vorschein, mit goldenen verschnörkelten Runen übersät. Doch blieb der wunderbare Runenmantel zusammengeknautscht liegen, als wäre er noch im Paket eingewickelt.

Leise winkte Cybelle Rangnir zu sich. Gemeinsam schafften sie es, den widerspenstigen Stoff aufzufalten. 

"Warum ist der so schwer?", ächzte Rangnir.

"Weil du nicht nur ihn bewegst, sondern auch sein Gegenstück im Kerkerturm", meinte Cybelle. Sie hatte es ihm doch schon dreimal erklärt!

Mit einem Ruck drehten sie den Runenmantel um. Rangnirs Mund klappte auf. Auf der Rückseite war keinerlei Stoff zu erkennen. Stattdessen funkelte und wirbelte etwas offensichtlich Magisches darin. Wie auf einer ruhigen Seeoberfläche glitzerte in der Mantelrückseite der Sternenhimmel. Doch als Cybelle und Rangnir sich darüber beugten, blickten ihnen nicht ihr eigenen Gesichter entgegen. Sondern das des kleinen Praeda, der über beide Backen hinaus strahlte.

"Wahrhaftig!", ertönte Corsarias Stimme mit einem leichten Echo dahinter. "Ein hadrisches Portal, gewebt in echten Stoff! Wunderbar!"

"Das ist kein echter Stoff", grinste Cybelle. "Das ist eigentlich nur Wasser. Danwars \gend{Wassermagier*innen}{Wassermagier} könnte Stoffe und Talismane aus einer Art festen Wassers herstellen. Verstehe jemand die Regeln der Magie. Ihren eigenen Körper zu Wasser machen konnten sie ja offenbar problemlos, aber andere Gegenstände mitnehmen war dann auf einmal ein Problem."  

Sie langte mit der linken Hand in das Portal hinein. Es fühlte sich glitschig an, als es über ihre Haut glitt. Die Luft auf der anderen Seite war wärmer. Eine Hand packte die ihre, Cybelle zog kräftig, und Minerma wirbelte durch das Portal, aus dem Kerkerturm auf die \textsc{Lacerta}. Sie unterdrückte ein Lachen über ihre wiedergefundene Freiheit. Dann machte sie sich daran, das Schiff fahrtüchtig zu machen.

Als nächste war Corsaria dran. "Ist dieses Portal dann auch nur Wasser?", fragte sie, sobald sie auf dem knarrenden Deck stand.

"Nee, hadrisch-magisch", meinte Cybelle nur. "Solche Portale bilden sich dort spontan im Schnee aus Ewigen Feuern. Irgendwie gelang es \gend{einer Wassermagier*in}{einem Wassermagier} und \gend{einer hadrischen Zauber*in}{einem hadrischen Zauberer}, hadrische Schneekristalle in ein Wassergewand einzuweben. Und die Runen halten es unnatürlich lange offen. Eine einzigartige Kollaboration zwischen zwei magischen Traditionen des Hadrischen Meeres. Kurzlebig, denn den Danwaren wurde anschließend rasch klar, dass durch solche Erfindungen vor allem mehr \gend{Hadrier*innen}{Hadrier} nach Danwar gelangen konnten, während die Danware mit Wassermagie auch ohne solche Tricks große Distanzen zu überwinden vermochten. So will es zumindest die Sage. Sie sind ein Vermögen wert."

"Wem hast du den denn beim Glücksspiel abgeknöpft?", fragte der kleine Praeda, welcher aufgeregt vor sich hin wippte. 

"Der Trägerin des größten Lavasteins von ganz Danwar. Aber nicht im Glücksspiel. Das ist eine Geschichte für ein andermal", meinte Cybelle.

Immer mehr Leute stürmten aus dem Portal und verteilten sich über Bord. Corsaria flüsterte Befehle, welche ebenso ruhig befolgt wurden. Der Anker wurde hochgezogen, die Segel gehisst, das Steuer befraut. Nicht nur Cybelles Crew entkam aus dem Kerkerturm, so manche stürzten sich ins Wasser und machten sich auf den Weg zurück zum sicheren Land.

Ein starker Nordwestwind bauschte die Segel. Die \textsc{Lacerta} nahm rasch Fahrt auf. Cybelle blickte argwöhnisch durch ihr Fernrohr, in der Hoffnung, dass es keine verfolgenden \gend{Seekrieger*innen}{Seekrieger}-Schiffe zeigte. Als sie sich dieser Hoffnung vergewissert hatte, zückte sie eine Flasche alten Rums und trat vorne zum Bug.

Cybelle beugte sich über den Bug und las die großen Buchstaben der \textsc{Lacerta}. 
"Wir brauchen keinen Namen, der uns nichts sagt". Cybelle ließ eine Flasche am Schiff zerschellen und proletete: "Ab heute seist du bekannt unter dem Namen \textsc{Sturmechse}! Unsere Nachfolgerin der vernichteten \textsc{Perlenechse}. Wir werden Großes mit dir erreichen." Bei nächstbester Gelegenheit würde sie die Segel mit dem Symbol der Schwarzen Kogge bemalen müssen.

Sie atmete tief durch, ließ sie ihren Blick über das prunkvolle Deck schweifen, zählte ihre eigene Crew durch ... und merkte überrascht, dass niemand fehlte. Im Gegenteil, es war jemand zu viel dabei.
Ein glatzköpfiger Bursche, zwischen drei und vier Jahrdutzend alt. Er starrte Cybelle unverwandt aus grauen Augen an, ein leichtes Lächeln auf den Lippen.

"Das ist Alfreck", meinte der kleine Praeda, der ihren Blick bemerkt hatte. "Er kann verdammt gut Karten spielen, das hat uns so einige erfreuliche Nacht im Gefängnis beschert. Ich habe meine letzten Socken an ihn verloren."

"Deswegen habe ich auch gar nicht erst gegen ihn gespielt", knurrte Minerma. "Er hat zu viel Glück!"

Das war offenbar laut genug gewesen, sodass Alfreck sie gehört hatte. Er schlenderte zum Steuer hoch und sprach mit hoher Stimme: "Glück habe ich tatsächlich vieles. Aber Kartenspiele sind nicht der Grund, warum ich beschlossen habe, noch ein bisschen länger in eurer Gesellschaft zu bleiben ... wenn ich mich denn haben wollt." 

Cybelle verschränkte ihre Arme und wartete auf weitere Worte.

"Zum einen wollte ich mich ganz außerordentlich bedanken bei meiner Befreierin. Eine wahrlich meisterhafte Leistung, die ihr hier abgezogen habt. Dieser Mantel ist ein unglaublich wunderbares Artefakt ..."

Schmeicheln konnte er. Cybelle mochte keine Schmeichler. Das schien Alfreck auch zu spüren, denn nun ließ er seine Schultern senken und änderte er seinen Ton. Schnarrend fuhr er fort: "Und, um ganz ehrlich zu sein, nur ein Verschwender lässt sich die Gelegenheit entgehen, auf einem Schiff von Seemann Emmetts Flotte die Hadrischen Meere zu überqueren. Erst recht, wenn die Alternative darin besteht, sich schifflos in Werftheim zu verstecken, verfolgt vom gesamten \gend{Seekrieger*innen}{Seekrieger}-Wachcorps. Was meint ihr, wie viele der armen befreiten \gend{Schlucker*innen}{Schlucker} die \gend{Seekrieger*innen}{Seekrieger} binnen zweier Tage bereits wieder eingeknastet haben?"

Nun war dieser Alfreck wenigstens ehrlich. Cybelle musterte ihn etwas eindringlicher. Er war drahtig, aber unter seiner dünnen Kleidung zeichneten sich auch Muskeln ab. Er konnte sich körperlich betätigen, das war klar.

"Ich werde Euch nicht zur Last fallen, im Gegenteil", versicherte Alfreck lächelnd. "Ich habe viel Erfahrung auf Schiffen, das liegt mir im Blut. Schon mein ... ne Großmutter war Kapitänin. Gemeinsam kommen wir weiter als alleine."

"Ich wäre ja nicht alleine", meinte Cybelle gelassen. "Wofür haben die \gend{Seekrieger*innen}{Seekrieger} dich schuften lassen?"

Alfreck zögerte einen Augenblick, ehe er antwortete: "Diebstahl. Irgendwie müssen wir ja über die Runden kommen. Geklaut von jenen, die die Absenz ihres Schmucks gar nicht gemerkt hätten."

"Offensichtlich haben sie es doch gemerkt", meinte Rangnir trocken. Er hatte sich noch immer nicht vom Konzept der Unantastbarkeit persönlichen Besitzes gelöst. Cybelle ignorierte ihn.

Erneut meldete sich der kleine Praeda: "Alfreck, du hast ja noch gar nicht von der Karte erzählt!"

Cybelle hob eine Augenbraue. "Spielkarte?"

Alfreck grinste und schüttelte seinen Kopf. "Schatzkarte. Wenn ihr mich Teil eurer Crew werden lässt, soll es euer Schaden nicht sein. Ich komme mit folgendem:"

Er zückte einige Streifen Pergament. Den obersten überreichte er Cybelle, die weiteren behielt er betont bei sich.

Leise raunte er ihr zu: "Die Karte eines Piratenkapitäns, welcher den größten Schatz des Nordens gefunden und aufgestockt haben will. Einen Hort, den Seekönig Varatan lange suchte, aber nie aufzuspüren vermochte. Das Versteck der Schwarzen Kogge."

Einen Herzschlag lang setzte Cybelles Herz aus. Natürlich war es eine Lüge, musste es eine Lüge sein. So viele Leute behaupteten schon, ein Versteck der Schwarzen Kogge gefunden zu haben. Der Großteil ihrer Plünderungen mussten auf der verfluchten Insel Narkon liegen, unerreichbar, verloren. Aber welche Piratenbande hatte nur ein Versteck? Schon alleine ein Bruchteil der errungenen Wertgegenstände und des Golds der Schwarzen Kogge könnte so viele \gend{Bewohner*innen}{Bewohner} der Nebelinseln in eine goldenere Zukunft führen.

"Woher hast du die Karte?", fragte Cybelle.

"Von einem garstigen Handelszwerg. Er würde mir nichts vormachen. Ich machte ihm klar, dass ich ihn aufsuchen und fertigmachen werde, sollte sich das Ding als Fälschung herausstellen."

Cybelle musterte die klaren Linien auf dem Schriftstück. Es war tatsächlich eine rudimentäre Karte, die einen Bereich des Hadrischen Meeres zeigte. Klippen und Inseln waren mit Punkten angedeutet, aber durchaus akkurat platziert. Entweder war das Werk hastig entstanden oder wollte Uneingeweihten nicht zu viel verraten. Die Buchstaben waren schwer zu lesen, da verschnörkelt. Eher zurückhaltend verfasst als hastig, demnach.

"R-A-T-T-E-N-Z-A...", las Cybelle langsam. Falls Alfreck davon überrascht war, dass sie die Zeichen lesen konnte, behielt er es für sich. Cybelle war der Rattenzahn wohlbekannt. Diese Klippenformation lag gefährlich weit im Westen. Maskierte Piraten und Riesenkraken waren keine Seltenheit dort. Und nun besagte diese Karte, dass dort ein Versteck Schätze ungeahnten Ausmaßes beherbergte.

Die hohe Stimme des Ausbrechers erklang erneut: "Genau. Dort beim Rattenzahn verbirgt sich der größte Schatz des gesamten Hadrischen Meeres."

Alfreck trat nahe an Cybelle heran, viel zu nahe. Er beugte sich vor, um den Rattenzahn auf der Karte zu zeigen, als hätte Cybelle ihn nicht schon längst erkannt. Seine Hand streifte wie zufällig ihren Rücken und blieb dort liegen. Cybelle klatschte ihm energisch mit der Hakenhand auf seine fleischige und zischte ihm zu: "Beim nächsten Mal verlierst du sie!" 

Alfreck trat einige Schritte zurück und rieb sich stumm die schmerzende Hand, seine Miene starr.

Cybelle knurrte, als sie überflog, was an erheblich leserlichen Notizen auf dem Pergament neben dem Versteck stand: "Es ist von einem Meisterschloss verdeckt. Sieben Schlüssel, drei davon magischer Natur, werden benötigt, um es zu öffnen. Das Versteck könnte gerade so gut auf dem Mond liegen, so wenig kommen wir dort hinein."

Alfreck fing sich wieder und lächelte kalt. "Das ist bei Weitem nicht alles. Zahlreiche Fallen schützen den Schatz. Ohne die detaillierten Beschreibungen der nächsten Seiten kommt ihr erst recht nicht unversehrt durch." Er schüttelte die restlichen Pergamente, die er für sich behalten hatte. Auf den ersten Blick handelte es sich dabei nicht um Karten, sondern um lange Texte mit einigen komplizierten Diagrammen.

"Doch um Schlösser bräuchten wir uns nur zu sorgen, wenn geschickte \gend{Zwerg*innen}{Zwerge} nicht Mittel erfunden hätten, um beinahe beliebige Schlösser zu öffnen. Die sieben Schlüssel können wir uns sparen. Und bezüglich der magischen Siegel kann ich versichern, dass ich eines sicher überwinden kann, und für ein zweites zumindest einen konkreten Lösungsansatz kenne."

Cybelle wusste durchaus, was für Wunderwerkzeuge \gend{Zwerg*innen}{Zwerge} produzieren konnten. "Ah, einen Casamatuc hast du? Ein wahres Meisterwerk zwergischer Mechanik. Ein ultimativer Dietrich."

"Ich besitze nicht direkt einen", druckste Alfreck herum. "Die sind hier im Norden schwer zu bekommen. Ach, was mache ich für Faxen, die sind überall schwer zu bekommen. Aber das soll uns nicht aufhalten. In Werftheim gibt es bekanntlich für genug Geld alles und ohne nichts. Und wenn man ein Dieb ist, gibt es in Werftheim schlicht alles."

Cybelle schüttelte entschieden ihren Kopf. "Auf keinen Fall wollen wir zurück nach Werftheim. Die \gend{Seekrieger*innen}{Seekrieger} werden uns abfangen, ehe wir die Insel auch nur sehen. Zum Glück kenne ich ..."

"Ich kenne jemanden!", fiel Rangnir ihr ins Wort. "Einen Silberzwerg, der für den richtigen Preis alles besorgen kann! Bestimmt auch einen Casamatuc."

"Einen Meisterdieb aus den Silberminen?", fragte Alfreck interessiert.

"Mitnichten", sagte Rangnir entrüstet. "Ein meisterhafter Hersteller aller möglichen Artefakte und Werkzeuge. Mein Dolch ist von ihm."

"Meiner ebenso.", sagte Cybelle. "Wir tragen beide je einen Dolch der Verbundenheit. Bilgram ist \textit{mein} Kontakt, Rangnir, ich hatte viel mehr mit ihm zu tun als du, du hast doch kaum einmal Werftheim verlassen."

"rAngnIr!", ertönte ein geisterhaftes Rauschen aus den Tiefen des Ozeans. Ein Rumoren rumpelte durchs gesamte Schiff. Wassertropfen wanderten den Bug hinauf, quollen durch die Ritzen des Decks und versammelten sich vor Cybelle in Form einer menschlichen Gestalt. Kirj'vana, der Meeresgeist, war zurückgekehrt. Mist, sie hatte Cybelle völlig verdrängt.

Kirj'vana ließ ihren wässrigen Zeigefinger über die Crew wandern. Selbst Alfreck hielt ausnahmsweise einen Kommentar zurück, als der Meeresgeist über das Deck wogte und um die verschiedenen Crewmitglieder waberte. Schließlich hielt sie vor dem bibbernden Rangnir an. "Du bist der Rangnir?"

Jener zitterte so stark, dass sein Nicken kaum zu erkennen war.

"Du hast der da drüber den Mantel mit Ruufs Knöpfen geschenkt?"

Erneut nickte Rangnir zitternd. "Sie hat ihn mir im Astzt-Spiel abgeknöpft."

Ihre Stimme wurde heller, als sie Cybelles Worte von vor einigen Wochen wiederholte: "Er ist ein Schiffsbauer aus Werftheim. Er hat vielleicht noch mehr Andenken an Ruuf. Wenn du mich an Land bringst, kann ich ihn zu dir holen."

"Und geholt habe ich ihn. Andere Artefakte von Ruuf hatte er leider nicht", meinte Cybelle. Sie musste die Kontrolle über diese Situation wieder an sich reißen. Dies war Kirj'vanas Element.

Sie ignorierend fragte Kirj'vana Rangnir: "Und woher hattest du die Knöpfe?"

Er schlotterte weiterhin zu sehr, was Cybelle nur zu gerne als Gelegenheit nahm, einzugreifen. "Er hat sie von einem Silberzwerg namens Bilgram erhalten. Wir sind jetzt auf dem Weg zu ihm."

An die Crew gewandt, hob Cybelle ihre Stimme und sprach: "Das ist Kirj'vana, ein mir wohlbekannter Meeresgeist. Fürchtet euch nicht vor ihr. Sie ist eine von uns, wenn sie denn will."

An Kirj'vana gewandt, fügte sie an. "Du kannst dich uns gerne anschließen auf dem Weg nach Silberhall. Sofern du dich an Bord halten kannst, während wir in Windeseile übers Meer segeln."

Falls Kirj'vana eine Lüge vermutete, zeigte sie es nicht. Ein letztes Mal wogte sie argwöhnisch um Rangnir, welcher Cybelle nur mit großen Augen ansah. Dann nickte sie kaum erkennbar. "Ein Flussgeist bin ich nicht. Das offene Meer ist meine Gewöhnung. An Bord bleiben kann ich. Das Ziel ist sinnvoll. Die Silberminen sind am südlichsten Zipfel Narkons. Ruuf muss dort nahe sein. Vielleicht sah dieser Bilgram ihn."

Erleichtert atmete Cybelle auf, als Kirj'vana zum Sturmmast hinüber wogte. Dort verlor sie langsam ihre Form, während Wasserperlen das Sturmsegel benetzten.

Rangnir setzte sich wieder hin und wippte leise vor sich hin, immer wieder nervöse Blicke zum Sturmsegel werfend, wo Kirj'vanas Gestalt nicht mehr erkennbar war. Corsaria wickelte ihn kurzerhand in eine warme Decke, bis er nicht mehr zitterte. 

"Das war dramatisch", merkte Alfreck an. "Kirschwana heißt die Nette?"

"Fast. Das 'sch' muss schwach sein, wie ein Hauch, mit ein kurzen Glottisschlag danach, und dann hast du die Aussprache."

"Kein 't'? Da gehört doch ein 't' rein, das fühle ich." 

"Ich weiß nicht, wie sehr ich deinem Gefühl trauen will."

Stille machte sich breit. Das Plätschern der Wellen an den Rumpf der \textsc{Sturmechse} ließ das Schiff leise Schaukeln. Die Mehrheit der Crew saß an Deck und hatte Mühe damit, die Augen offenzuhalten. Manche waren bereits unter Deck gegangen, um die besten Hängematten für sich zu erringen. Auch Cybelle konnte ein Gähnen nur schwer unterdrücken. Es war längst Zeit, schlafen zu gehen.

Doch noch war es Zeit für eine Rede. Cybelle wandte sich ihrer Crew zu und sprach laut: "Heute sind wir den Ketten der Werften und den Gittern der Seekrieger entronnen. Diese Schmach werden wir nicht auf uns sitzen lassen und nicht vergessen. Das Volk der Nebelinseln bedarf noch immer unser Unterstützung. Doch wurde uns in dieser Begegnung etwas geschenkt, was noch wertvoller ist als einige erbeutete Zölle tätowierter Händlerinnen. Eine Schatzkarte zu einem der größten Schätze des Meeres. Eine Karte zum Versteck der Schwarzen Kogge!"

Jubel ertönte. Und Cybelle wandte sich ab, um die Nachtschicht einzuteilen.





\newpage
\section{Bilgram, der Silberzwerg}

Einen einzelnen Großhafen gab es am südlichen, nicht verfluchten Teil von Silberland. Natürlich war dort alles voller \gend{Hafenmeister*innen}{Hafenmeister}, die womöglich bereits durch Falken oder Runenstiften erfahren hatten vom Diebstahl von Emmetts Schiff. Für \textsc{Sturmechse} kam dort anzulegen nicht infrage.

Als nächstes versuchte Cybelle es am südlichsten Zipfel von Silberland. Auf einer weit ins Meer hinausragenden Landzunge lag dort Hallworts Ruhe, das Grab des Zwergenfürsten Hallwort, dessen Reise in den Norden die Silbermine und die Gemeinschaft der \gend{Silberzwerg*innen}{Silberzwerge} vor über 60 Jahren begründet hatte. Doch waren aktuell allerlei Wachen stationiert, mit leuchtenden Fackeln und zahlreichen Nachtpatrouillen. 

Anscheinend hatte jemand kürzlich die Grabruhe gestört, Hallworts Ruhe aufgebuddelt und dessen sicherlich wertvolle Grabbeigaben geklaut. Alfreck machte sich stark dafür, weit weg von hier zu ankern. Cybelle stimmte seiner weisen Vorsicht zu.

Letzten Endes ankerte die \textsc{Sturmechse} im Dunkel der Nacht, einige Meilen weiter nördlich, an der Bronzeküste.

Cybelle und Rangnir waren die einzigen, die per Ruderboot an Land gingen. Kirj'vana spazierte ihnen ganz gemütlich über die Wasseroberfläche nach und wartete als Silhouette am Meeresufer. Alfreck blieb mit der restlichen Crew an Bord der \textsc{Sturmechse} und hielt Ausschau nach verfolgenden \gend{Seekrieger*innen}{Seekrieger}-Schiffen. Oder er hatte einfach nicht so sehr Lust, sich die Beine zu vertreten, und zockte stattdessen lieber die restlichen Besatzungsmitglieder in einigen Runden Astzt königlich ab. Cybelle schüttelte ihren Kopf. Sie sollten es doch besser wissen, als mit einem so offensichtlichen Falschspieler zusammenzusitzen.

Cybelle stapfte durchs seichte Wasser und zog mit Rangnir das Ruderboot ganz ans Land. Die Bronzeküste war mit steilen Klippen übersät, doch hatte sie die \textsc{Sturmechse} wohlweislich vor einem kleinen Sandstrand platziert. 

Sie griff mit ihrer linken, gesunden Hand in den Sand und ließ die rauen Körner durch die Finger rieseln, bis nur noch einige leere Muschelschalen übrig waren, die sich angenehm aneinander reiben ließen. Das Rauschen der Wellen entspannte sie. Dies waren aufregende Tage gewesen, doch nun war es geglückt. Sie waren befreit, ihre Crew und sie wieder vereint, und sie hatten ein neues aufregendes Ziel: Das Versteck der Schwarzen Kogge! Wenn Bilgram ihr nur einen Casamatuc verschaffen konnte. Und diesen Meeresgeist wieder loswerden musste sie auch noch. Aber auch hier konnte Bilgram vielleicht aushelfen.

Cybelle und Rangnir zogen natürlich nicht bis zum großen Eingangstor in die Silbermine. Stattdessen folgten sie einem gewundenen Pfad den Silberberg hoch. Zweimal mussten sie Trupps von \gend{Silberzwerg*innen}{Silberzwergen} ausweichen, die allerdings zu sehr in Konversationen miteinander vertieft waren, um die beiden \gend{Freibeuter*innen}{Freibeuter} zu erspähen. 

Schon einige Jahre lang war Cybelle nicht mehr hier gewesen, doch kam ihr alles noch bekannt genug vor. Die letzten paar hundert Meter mussten sie einen Berghang hochkraxeln, ohne dass irgendein Pfad den Weg gezeigt hätte. Cybelle kannte die Felsformation, die sie suchte. 

Sie drehte sich um und genoss den Ausblick auf die Bronzeküste hinter ihr. Rangnir erreichte sie keuchend. Lächelnd trat Cybelle einen Schritt zurück ... direkt in den Felsen hinein! Sie verspürte ein leichtes Kitzeln, als ginge man durch einen dünnen Vorhang aus kaltem Wind. Und dann: Dunkelheit! Der Zauber verbarg nicht nur die Zwergentür vor neugierigen Blicken, er ließ auch kein Licht hindurch.


"Ein Tarnzauber! Tut das weh?", fragte Rangnir ängstlich.
      
"Es kitzelt ein bisschen", meinte Cybelle. "Aber ansonsten fühle ich nichts Besonderes."

"Uääää", machte Rangnir, während er selbst durch den Tarnzauber trat.

Cybelle lachte. "\gend{Zwerg*innentüren}{Zwergentüren} zu unterirdischen Tunneln sind oft gut versteckt oder mit solchen Runenillusionen getarnt. Wir können von Glück sprechen, dass diese hier immer noch offen war, sonst hätten wir Bilgram anderswie kontaktieren müssen."

Nach ein paar Schritten in der absoluten Dunkelheit glomm eine grünliche Runenzeichnung unter Cybelles Fuß auf, breitete ihr Licht bis auf schlichte Runenkreis an den Wänden und tauchte den engen Gang in ein schwummriges Licht. Die Wände waren von den Runen abgesehen völlig glatt geschliffen, ohne viel Firlefanz, wie man es in den Gängen von Cavern erwarten würde. Die glühenden Runen ließen Rangnirs Gesicht gruselig aufleuchten. Dieser hatte sich wieder gefangen, schlug anerkennend an die feste Wand und meinte schlicht: 

"Toll. Ein bisschen Licht gibt es ja auch noch! Ein Hoch darauf, dass die \gend{Silberzwerg*innen}{Silberzwerge} nicht nur an Leute mit Dunkelsicht denken."

Gemeinsam schritten sie durch den dunklen Gang. Wo immer sie hintraten, glommen weitere dimme Runenzeichnungen an der Wand auf. Cybelle und Rangnir stießen sich einige Male die Köpfe an Vorhängen der niedrigen Decke. Hin und wieder kamen sie an Abzweigungen vorbei. Cybelle kam nicht umhin, sich zu wundern, wie viele Abzweigungen sich wohl hinter weiteren Tarnzaubern versteckten.

"Links, rechts, wieder rechts, dann ein letztes Mal links...", murmelte Cybelle leise, in der Hoffnung selbstsicherer zu wirken, als sie war. Hier hatte sie sich vor Jahren schon einmal verirrt.

Doch nicht diesmal.

Bald schon erreichten sie eine kleine quadratische Hinterkammer. An den Wänden standen klobige Steinsitze. Die Raummitte war eingenommen von einem Wasserbecken, in welches alle paar Sekunden ein Tropfen aus einem Loch in der Decke hinein platschte. Sowohl in den Sitzen als auch im Becken hätte Cybelle gerade knapp Platz gehabt. 

Rangnir ging in die Hocke und prüfte, ob es zwischen dem Becken und dem Boden eine Fuge gab, konnte aber keine finden. Offensichtlich wurde das Becken als Ganzes bei der Schaffung dieses Raums angelegt, und mit den Sitzen war es gleich – erstaunlich!

Mit ausgestreckter Hand trat Rangnir näher ans Becken. "Fass es nicht an", rief Cybelle rasch, "Es würde die Interaktion weiterleiten. Wir wollen offiziell nicht hier gewesen sein." 

Stattdessen ging sie zur Tür und klopfte eine bestimmte Klangfolge daran. Das wiederholte sie jede Viertelstunde, bestimmt mehrere Stunden lang, bis schließlich endlich auf der anderen Seite ein Knarren und Knirschen von der Öffnung der Steintür zeugte.

"Endlich!", seufzte Cybelle. Rangnir, welcher zunächst ihren ganzen Proviant an trockenem Brot verspeist hatte und danach schnarchend an eine Wand gesunken war, schreckte auf.

Die Tür schwang auf und gab den Blick frei auf einen großen Arbeitsraum. Jede mögliche Fläche war über und über überstellt mit Holz- und Metallstücken in verschiedensten Formen. An den Wänden hingen Skizzensammlungen und mit Kreidezeichen überstellte Steintafeln. Zahlreiche Schaukästen enthielten verschiedenste Waffen, die meisten mit aufwändigen Verzierungen und leuchtenden Farben.

Und mitten im Raum, an einem großen Hebel, welcher soeben die Öffnung der Tür in Gang gesetzt hatte, stand ein Silberzwerg mit einem braunen Bart und blickte sie argwöhnisch an.

Cybelle und Rangnir zückten ihre Dolche der Verbundenheit und schlugen sie zusammen. Das beruhigende Summen der vibrierenden Dolche übertönte den Lärm eines ratternden Konstrukts am anderen Raumende, welches grüne Flammen spuckte. Da tauchte ein breites Lächeln auf dem Gesicht des Zwergs auf.

"Cybelle. Wie lange ist es her. Und ... und La ... Ri ..." 

Er sah betreten zu Boden.

"Rangnir", sagte Rangnir, "Werter Bilgram, ihr wart mal zu Besuch in Werftheim und habt meine Vorsteherin schwer beeindruckt mit euren Silberkonstrukten. Und mich."

Wenn es irgendwie möglich war, sah der Zwerg nur noch betretener aus.

Lange fackelte Cybelle nicht herum. "Werter Bilgram, verzeih die späte Störung, doch wussten wir nicht, wem wir vertrauen können. Wir sind hier ..."

Bilgram winkte sie ab und hinein in die gute Stube. Rasch schob er einige Skizzenstapel und lose Schrauben von zwei Sitzkissen und gebot den beiden \gend{Besucher*innen}{Besuchern}, sich zu setzen. Er füllte Kräutertee aus einem meterhohen Rohr in Metallbecher, servierte sie ihnen auf einem schmierigen Tablett und blieb vor ihnen stehen, während er sprach: "Ah, Cybelle. Als ich die Nachricht von Werftheim vernahm, ahnte ich, dass du dahinter stecktest. Dein magischer Mantel ist einmalig. Ich hoffe doch stark, die \gend{Seekrieger*innen}{Seekrieger} können den  nicht auf mich zurückführen. Es dauerte lange, die kaputten Runen nachzuziehen. Und verloren habt ihr ihn auch noch."

Cybelle schüttelte ihren Kopf: "So wie ich dich kenne, findest du den doch schon bald wieder im Besitz irgendeines Handelszwergs, und kannst ihn zurückkriegen und für einen ebenso horrenden Preis wieder verleihen."

"Spannend. War gratis für dich so ein horrender Preis?"

"Spannend. Waren meine Gegenleistungen für dich ebenso gratis? Oh, wie lange meine Hand die deine waschen musste dafür."

Rangnir gähnte. "Ja, ja, ihr beide kennt euch lange, wir verstehen's. Kommt zur Sache."

"Mein lieber ... Rangnir, lass mich dich nicht langweilen", sprach Bilgram. "Ein Schiffsbauer bist du? Würde es dir mehr behagen, dich in meiner Werkstatt umsehen? So manches Stück hier könnte für dich spannend sein."



Nun ließ auch Cybelle ihren Blick über Bilgrams Werkstatt schweifen. Es hatte sich einiges verändert seit ihrem letzten Besuch hier. Sie stupste mit ihrem Fuß einen Bogen an, der achtlos am Boden lag.

"Ich kann mir nichts hier drinnen leisten", murmelte Rangnir. "So viele edle Handwerksmaterialien. So viel Prunk. Du musst im Gold baden können."

Bilgram schüttelte nur traurig seinen Kopf. "Dass ich damit arbeite, heißt nicht, dass ich es danach behalten kann. Ich komme knapp über die Runden. Und ich stelle das her, was man von mir verlangt. Waffen, größtenteils. Alle wollen Waffen, darin liegt das große Geld. Warum nicht mal ein mechanischer Arm, der eine Kinderkrippe periodisch anstößt? Warum nicht eine Statue an die Liebe? Nein, nichts als Werkzeuge des Tötens."

"Wir brauchen kein Werkzeug des Todes. Wir sind hier für einen Casamatuc", kam Cybelle zur Sache.

"Schwierig. Gefährlich. Teuer", schoss Bilgram zurück. Er kramte eine Zeit lang in einer Schublade herum und präsentierte dann ein kleines,
unscheinbares Stück Metall mit einigen Scharnieren und Drähten dran. Einen waschechten Casamatuc! Cybelles Augen leuchteten auf. Bilgram platzierte den Casamatuc hingegen betont ruhig auf den Tisch zwischen den beiden. Die Verhandlungen waren eröffnet.

Rangnir gähnte. Cybelle appellierte an Bilgrams Ego, seine Erfinderehre, seinen Abenteuerdurst. Vergeblich. Erst als sie von einer Schatzkarte zum Versteck des größten Schatzes der Nebelinseln zu sprechen kam, zuckte Bilgrams Blick auf.

"Könnt ihr für die Authentizität der Karte bürgen?", fragte Bilgram. "Woher habt ihr sie?"

"Ein Dieb namens Alfreck hat sie uns zur Verfügung gestellt, dafür dass wir ihn aus den Zwangswerften gerettet hatten."

"Äußerst vertrauenswürdig."

"Natürlich nicht, aber die Karte zeigt mir wohlbekannte Klippen äußerst akkurat. \gend{Ihre Erschaffer*in}{Ihr Erschaffer} war viel auf Hoher See. Alfreck meinte auch, seine Großmutter sei bereits Kapitänin gewesen. Wenn dieser Glatzkopf nur mal die Wahrheit sagt ... er hat eine Zunge aus Silber, die meisten meiner Leute scheinen ihn zu lieben, aber er gibt mir eine Gänsehaut. Wenn es dir nicht behagt, mit einem Dieb zusammenzuarbeiten, hätte ich auch nichts dagegen, ihm einfach die Karte abzuknöpfen."

Bilgram zuckte betont ruhig mit den Schultern. "Kommt auf den Dieb an. Ihr habt mir demonstriert, dass \gend{Pirat*innen}{Piraten} auch angenehme Gesellschaft sein und das Herz am rechten Fleck haben können. Was hat dieser Dieb denn gestohlen?"

"Schmuck, sagt er. Er ging sparsam mit Details um."

"Ein glatzköpfiger Schmuckdieb, der in Werftheim gefangen war, und eine Schatzkarte zum größten Schatz der Nebelinseln haben will", wiederholte Bilgram, seinen Bart nachdenklich streichend. Dann schlug er auf den Tisch. "Donnerwetter! Ich unterstütze euch, Cybelle. Mein Casamatuc sei euch jedoch nur geliehen. Ich werde ihn mit mir führen." Ein schiefes Grinsen umspielte seinen Mund, als er sprach: "Ich komme mit euch!"

Rangnir schreckte von seinem Schnarchen hoch. "Wie, was, Donnerwetter, du kommst mit? Hat dich das Abenteuerfieber gepackt? Willst du dir einen größeren Teil des Schatzes einsacken?"

"So was in der Art", lächelte Bilgram mysteriös. Dann hastete er los, seine Sachen zusammenzupacken. Er verteilte auch recht frei verstärkte Waffen und Ausrüstungsgegenstände auf seine \gend{Besucher*innen}{Besucher}.

Cybelles Blick fiel auf einen edlen Dolch, der in einem Glaskasten präsentiert war. Er war eine meisterliche Arbeit, die Klinge sah ebenso hart wie scharf aus. Das beeindruckendste jedoch war das Griffstück: eine goldene Schlange wand sich um den Griff, jede Schuppe war einzeln herausgearbeitet. Als Knauf fungierte der Kopf der Schlange, zwei Rubine leuchteten als glühend rote Augen, die bösartig funkelten. Eine gespaltene Zunge ragte aus dem Maul. Sie leckte sich die Lippen.

"Denk nicht mal dran", lachte Bilgram, "Ich habe die Hoffnung noch nicht aufgegeben, dass derjenige, der sie in Auftrag gab, sie eines Tages abholen kommt. Und mir die versprochene Belohnung herausrückt. Das kannst du dir nicht leisten. Außerdem ist er ein viel zu schönes Stück, um eine gute Waffe zu sein."

Doch hatte Bilgrams Auftraggeber den Schlangendolch wohl auch nicht wirklich hatte als Waffe nutzen wollen, mehr als Ausstellungsstück. Nun gut, bezahlen hatte er sie offensichtlich auch nicht wollen, und nun saß Bilgram auf einer nutzlosen Giftwaffe. 

Kurzerhand öffnete Bilgram den Glaskasten mit dem Dolch und griff selbst danach. Cybelle lachte anerkennend.

Während Cybelle und Rangnir sich abwandten, um einen farbenprächtigen Helm begutachteten, öffnete Bilgram das kleine braune Schäftlein, auf welchem der Glaskasten stand. Zwei Flaschen standen darin: Eine bauchige Glasflasche, in der eine dickflüssige schwarze Flüssigkeit schwappte. Daneben eine unscheinbare Phiole mit klarem Gegengift.

\trennzeichen

Bevor sie sich wieder auf den Weg nach draußen machen würde, mussten die erschöpften \gend{Besucher*innen}{Besucher} schlafen und sich etwas erholen, verlangte Bilgram. Er räumte einige behelfsmäßige Matten beiseite, bereite ihnen einen leckeren Schlaftrunk mit Cocoa und Honignusssaft und bedeckte sie mit kunstvollen Decken, die trotz ihrer Dünne gut wärmten.

Schon bald schallte Cybelles und Rangnirs Schnarchen durch die überstellte Werkstatt. Und Bilgram trat zu einem Pult und zog an einer Reihe von Schnüren, um eine Botschaft an die Silbergarde zu übermitteln. Einen Verdacht und einen Befehl.
 
Als das Trio später erfrischt durch die Geheimtür schlüpfte, blieb Bilgram einige Schritte hinter den anderen zurück. Ganz ungeniert lehnte er sich ans Steinbecken, in welchem Wasser leise vor sich hin plätscherte. Den Rücken zum Becken gekehrt, tauchte er seine Finger hinein und berührte das Wasser in einer ganz bestimmten Abfolge. Die Geheimcodes der Silberzwerge waren äußerst effizient: Einige Augenblicke genügten, damit Bilgram ausdrücken konnte, was er wollte. Ein Trio von Tropfen stürzte von der Decke und pochte auf die Wasseroberfläche. Ein Zeichen, dass  Bilgrams stumme Botschaft an die Silbergarde angekommen war. Sie wussten, dass die drei nun aufbrachen. 

Cybelle und Rangnir war nichts aufgefallen. "Kommst du nun oder nicht?", erklang Cybelles mürrische Stimme von weiter vorne. 

"Bin gleich da!", rief Bilgram. Er trocknete sich die Hände an seinem Mantel und rannte den beiden nach.

\trennzeichen

Die Morgensonne stand bereits hoch am Himmel, als das Trio durch die verborgene \gend{Zwerg*innen}{Zwergen}tür nach draußen schlüpfte. Erneut machte Rangnir "Uääää", als die kitzelnde Illusion über seine Haut glitt. Cybelle räkelte und streckte sich, ohne dass ihre Hände an eine zu tiefe Decke gestoßen wären. Bilgram setzte seinen viel zu großen Rucksack ab und verschnaufte.

Da brodelte eine magische Wassermasse hinter einem Felsen hervor und stakste auf die drei zu. Wo immer ihre Füße auf den steinigen Untergrund Silberlands trafen, dampfte und zischte es. Kirj'vana wagte sich tatsächlich auf das Land. Sie blubberte und zuckte, als würde die Bewegung sie schmerzen, aber der Meeresgeist wagte sich aufs Land, aufs verfluchte Silberland, wo vor vielen Jahrzehnten ihr Geliebter gefangen genommen war und vermutlich längst tot war.

Cybelle knirschte mit den Zähnen. Damit hatte sie nicht gerechnet. Nun war es zu spät, Bilgram anzuweisen.

Bedrohlich richtete sich Kirj'vana vor den dreien in die Höhe. "Bist du Bilgram?", fragte sie Bilgram. Bilgram nickte. Der Silberzwerg sah überrascht aus, aber nicht ängstlich. Er starrte Kirj'vana vielmehr beeindruckt an.

"Woher hast du Ruufs Knöpfe?", blubberte Kirj'vana.

"Welche Knöpfe?", wollte Bilgram wissen.

Und Rangnir fiel zitternd auf die Knie. "Es tut mir leid, o holder Meeresgeist Kirj'vana, ich habe die Knöpfe nicht von Bilgram. Ich habe gelogen."

Da fasste Cybelle sich ein Herz: "Nein, ich war es, die gelogen hat. Ich wollte hierher, ich wollte dich am Riemen halten. Ich habe meine eigenen Prioritäten. Und näher an Ruuf als nach Silberland kommst du ohnehin nicht. Aber die Knöpfe stammen anderswoher."

Kirj'vana schüttelte ihren Kopf, sodass Wasser überallhin spritzte. Sie sah nicht enttäuscht aus. Vermutlich hatte sie es schon längst gewusst.










\newpage
\section{Ruuf, der Seekrieger}



















\newpage
\section{Ken Dorr, der Dieb}


























\newpage
\section{Epiloge}






















